% ─── AUTO-GENERATED by Tosh CommandLatexEmitter ───────────────────────
% Do not edit by hand. Regenerate with: dotnet run --project src/Tosh.Cli -- --export-command-manifest --latex


\section{CLR}

\subsection{\texttt{alloc}}
\label{ref:alloc}
\icmd{alloc}

\begin{cmdbox}{alloc}
Allocates a native unmanaged buffer by byte size or interop type.

\textit{Aliases:} \code{native-alloc}

\begin{signature}
alloc <bytes | type-name>
\end{signature}

\end{cmdbox}

\subsection{\texttt{call}}
\label{ref:call}
\icmd{call}

\begin{cmdbox}{call}
Invokes an instance or static CLR method.

\begin{signature}
call <method-name> [args...] or call <type-name> <method-name> [args...]
\end{signature}

\end{cmdbox}

\subsection{\texttt{call-method}}
\label{ref:call-method}
\icmd{call-method}

\begin{cmdbox}{call-method}
Invokes a method by dynamic name.

\begin{signature}
call-method [object] <method-name> [args...]
\end{signature}

\begin{lstlisting}
echo hello | call-method ToUpper
call-method $obj MethodName arg1
\end{lstlisting}

\end{cmdbox}

\subsection{\texttt{cast}}
\label{ref:cast}
\icmd{cast}

\begin{cmdbox}{cast}
Casts pipeline values to a CLR type, including constructed generic collection types.

\begin{signature}
cast <type> [value ...]
\end{signature}

\begin{lstlisting}
echo [1, 2, 3] | cast list<int>
echo 42 | cast string
\end{lstlisting}

\begin{notebox}
Cast converts to CLR target types, including constructed generic collection types like `list<int>`.
\end{notebox}

\end{cmdbox}

\subsection{\texttt{clone}}
\label{ref:clone}
\icmd{clone}

\begin{cmdbox}{clone}
Creates a shallow copy of an object.

\begin{signature}
clone [object]
\end{signature}

\begin{lstlisting}
$obj | clone
clone $obj
\end{lstlisting}

\end{cmdbox}

\subsection{\texttt{constructors}}
\label{ref:constructors}
\icmd{constructors}

\begin{cmdbox}{constructors}
Lists constructors for CLR types, ToSh classes, and shell collection types.

\begin{signature}
constructors [type ...]
\end{signature}

\begin{lstlisting}
constructors System.String | first 5
constructors list<int>
constructors dict<string, int>
\end{lstlisting}

\begin{notebox}
Constructors works with CLR types, ToSh named types, and shell collection types like `list<int>` and `dict<string, int>`.
\end{notebox}

\end{cmdbox}

\subsection{\texttt{del-prop}}
\label{ref:del-prop}
\icmd{del-prop}

\begin{cmdbox}{del-prop}
Removes a property from a dynamic record.

\begin{signature}
del-prop [object] <name>
\end{signature}

\begin{lstlisting}
$obj | del-prop Name
del-prop $obj Name
\end{lstlisting}

\end{cmdbox}

\subsection{\texttt{describe-type}}
\label{ref:describe-type}
\icmd{describe-type}

\begin{cmdbox}{describe-type}
Describes CLR types, ToSh named types, or shell collection types.

\begin{signature}
describe-type [type ...]
\end{signature}

\end{cmdbox}

\subsection{\texttt{get-methods}}
\label{ref:get-methods}
\icmd{get-methods}

\begin{cmdbox}{get-methods}
Lists method names for an object or ToSh class.

\begin{signature}
get-methods [object]
\end{signature}

\begin{lstlisting}
$obj | get-methods
get-methods $obj
\end{lstlisting}

\end{cmdbox}

\subsection{\texttt{get-prop}}
\label{ref:get-prop}
\icmd{get-prop}

\begin{cmdbox}{get-prop}
Gets a property value by dynamic name.

\begin{signature}
get-prop [object] <name>
\end{signature}

\begin{lstlisting}
$obj | get-prop $propName
get-prop $obj Name
\end{lstlisting}

\end{cmdbox}

\subsection{\texttt{get-props}}
\label{ref:get-props}
\icmd{get-props}

\begin{cmdbox}{get-props}
Lists property names for an object.

\begin{signature}
get-props [object]
\end{signature}

\begin{lstlisting}
$obj | get-props
get-props $obj
\end{lstlisting}

\end{cmdbox}

\subsection{\texttt{has-method}}
\label{ref:has-method}
\icmd{has-method}

\begin{cmdbox}{has-method}
Checks whether an object or ToSh class has the named method.

\begin{signature}
has-method [object] <name>
\end{signature}

\begin{lstlisting}
$obj | has-method ToString
has-method $obj ToString
\end{lstlisting}

\end{cmdbox}

\subsection{\texttt{has-prop}}
\label{ref:has-prop}
\icmd{has-prop}

\begin{cmdbox}{has-prop}
Checks whether an object has the named property.

\begin{signature}
has-prop [object] <name>
\end{signature}

\begin{lstlisting}
$obj | has-prop Name
has-prop $obj Name
\end{lstlisting}

\end{cmdbox}

\subsection{\texttt{load-assembly}}
\label{ref:load-assembly}
\icmd{load-assembly}

\begin{cmdbox}{load-assembly}
Loads a .NET assembly from disk into the current process.

\begin{signature}
load-assembly <path> [path...]
\end{signature}

\end{cmdbox}

\subsection{\texttt{members}}
\label{ref:members}
\icmd{members}

\begin{cmdbox}{members}
Lists public members for CLR types or pipeline objects.

\begin{signature}
members [type ...]
\end{signature}

\begin{lstlisting}
members string
DateTime.Now | members
\end{lstlisting}

\end{cmdbox}

\subsection{\texttt{methods}}
\label{ref:methods}
\icmd{methods}

\begin{cmdbox}{methods}
Lists public methods for CLR types or pipeline objects.

\begin{signature}
methods [type ...]
\end{signature}

\end{cmdbox}

\subsection{\texttt{native-free}}
\label{ref:native-free}
\icmd{native-free}

\begin{cmdbox}{native-free}
Frees one or more native buffers allocated by native-alloc.

\begin{signature}
native-free [buffer ...]
\end{signature}

\end{cmdbox}

\subsection{\texttt{native-read}}
\label{ref:native-read}
\icmd{native-read}

\begin{cmdbox}{native-read}
Reads a C string, byte range, or native scalar/struct-layout value from native memory.

\textit{Aliases:} \code{read-buffer}

\begin{signature}
native-read <cstring|bytes|type-name> [buffer|pointer] [length] [offset]
\end{signature}

\end{cmdbox}

\subsection{\texttt{native-write}}
\label{ref:native-write}
\icmd{native-write}

\begin{cmdbox}{native-write}
Writes a C string, byte sequence, or struct-layout value into native memory.

\textit{Aliases:} \code{write-buffer}

\begin{signature}
native-write <buffer|pointer> <value> [offset]
\end{signature}

\end{cmdbox}

\subsection{\texttt{new}}
\label{ref:new}
\icmd{new}

\begin{cmdbox}{new}
Constructs a new CLR object, ToSh named type, or shell collection.

\begin{signature}
new <type-name> [ctor-args...]
\end{signature}

\begin{lstlisting}
var rng = new System.Random()
var items = new list<String>("one", "two")
new System.Text.StringBuilder("hello").Append(" world").ToString()
\end{lstlisting}

\begin{notebox}
Tosh supports both the legacy `new <Type> ...` command form and the newer C\#-style `new Type(...)` expression syntax. Shell collection types also support generic construction like `new list<String>(...)`.
\end{notebox}

\end{cmdbox}

\subsection{\texttt{offset-of}}
\label{ref:offset-of}
\icmd{offset-of}

\begin{cmdbox}{offset-of}
Returns the unmanaged field offset for a sequential or explicit-layout struct.

\textit{Aliases:} \code{native-offsetof}

\begin{signature}
offset-of <type-name>[.<field-name>] [field-name]
\end{signature}

\end{cmdbox}

\subsection{\texttt{set-prop}}
\label{ref:set-prop}
\icmd{set-prop}

\begin{cmdbox}{set-prop}
Sets or adds a property on a dynamic record.

\begin{signature}
set-prop [object] <name> <value>
\end{signature}

\begin{lstlisting}
$obj | set-prop Name "value"
set-prop $obj Name "value"
\end{lstlisting}

\end{cmdbox}

\subsection{\texttt{size-of}}
\label{ref:size-of}
\icmd{size-of}

\begin{cmdbox}{size-of}
Returns the unmanaged size of a supported native interop type.

\textit{Aliases:} \code{native-sizeof}

\begin{signature}
size-of <type-name> [type-name ...]
\end{signature}

\end{cmdbox}

\subsection{\texttt{type-of}}
\label{ref:type-of}
\icmd{type-of}

\begin{cmdbox}{type-of}
Returns the CLR type or ToSh class type for each input object.

\begin{signature}
type-of [value...]
\end{signature}

\end{cmdbox}

\subsection{\texttt{types}}
\label{ref:types}
\icmd{types}

\begin{cmdbox}{types}
Lists available CLR and ToSh shell types.

\begin{signature}
types [-a] [filter]
\end{signature}

\begin{lstlisting}
types System.String
types list
types map | where _.Namespace == ToSh
\end{lstlisting}

\begin{notebox}
Types searches both CLR types and ToSh shell types like `list`, `array`, `dict`, `table`, and `tuple`.
\end{notebox}

\end{cmdbox}


\section{Data}

\subsection{\texttt{as-file}}
\label{ref:as-file}
\icmd{as-file}

\begin{cmdbox}{as-file}
Materializes values into a temporary file and returns it as a file system entry.

\begin{signature}
as-file [text|json|csv] [value ...]
\end{signature}

\begin{tabular}{@{}>{\ttfamily\small\raggedright\arraybackslash}p{5cm} L{9cm}@{}}
\toprule
\textbf{Argument} & \textbf{Description} \\
\midrule
format & Output format: text, json, or csv. Defaults to text. (optional) \\
value & Values to materialize. Can also come from pipeline. (optional) \\
\bottomrule
\end{tabular}

\textbf{Pipeline input:} scalar, list, table
 --- Materializes pipeline values into a temporary file.

\textbf{Output:} Returns a FileSystemEntry for the created temporary file.

\begin{lstlisting}
ls | as-file  # Materialize pipeline to temp file
as-file json $data  # JSON format
\end{lstlisting}

\begin{notebox}
As-file materializes pipeline values into a temporary file and returns a file object you can pass to external executables.
\end{notebox}

\end{cmdbox}

\subsection{\texttt{from}}
\label{ref:from}
\icmd{from}

\begin{cmdbox}{from}
Parses structured text into objects.

\begin{signature}
from <format> [options] [text]
\end{signature}

\begin{lstlisting}
echo "{\"name\":\"toast\"}" | from json
curl https://example/api | from json | flatten
cat data.toml | from toml
cat data.csv | from csv
\end{lstlisting}

\begin{notebox}
The `from` and `to` commands convert between text formats (json, csv, tsv, xml, toml) and CLR objects. Parsed values stay as CLR objects until you explicitly flatten them.
\end{notebox}

\end{cmdbox}

\subsection{\texttt{hash}}
\label{ref:hash}
\icmd{hash}

\begin{cmdbox}{hash}
Computes hashes for text input or files.

\begin{signature}
hash [algorithm] [path ...]
\end{signature}

\end{cmdbox}

\subsection{\texttt{parse}}
\label{ref:parse}
\icmd{parse}

\begin{cmdbox}{parse}
Parses text input with a regular expression into shell record objects.

\begin{signature}
parse [-a] [-i] [-m] [-s] [-x] [--explicit-capture] <pattern|regex> \\
\quad [text ...]
\end{signature}

\begin{lstlisting}
ping -c 3 localhost | parse "time=(?<time_ms>[0-9.]+) ms"
echo "PID=42" | parse "PID=(?<Pid>[0-9]+)"
echo "first\nsecond" | parse -am "^(?<Value>\\w+)$" | get Value
\end{lstlisting}

\begin{notebox}
Parse and match use .NET regular expressions, including named groups and inline modifiers like `(?im)`.
\end{notebox}

\end{cmdbox}

\subsection{\texttt{to}}
\label{ref:to}
\icmd{to}

\begin{cmdbox}{to}
Serializes objects into structured text.

\begin{signature}
to <format> [options]
\end{signature}

\begin{lstlisting}
ls | to json
ls | to csv
ls | to toml
ls | to json --compact
\end{lstlisting}

\begin{notebox}
The `from` and `to` commands convert between text formats (json, csv, tsv, xml, toml) and CLR objects. Parsed values stay as CLR objects until you explicitly flatten them.
\end{notebox}

\end{cmdbox}


\section{Filesystem}

\subsection{\texttt{append-file}}
\label{ref:append-file}
\icmd{append-file}

\begin{cmdbox}{append-file}
Appends plain text to a file, creating it when needed.

\begin{signature}
append-file <path> [value...]
\end{signature}

\begin{tabular}{@{}>{\ttfamily\small\raggedright\arraybackslash}p{5cm} L{9cm}@{}}
\toprule
\textbf{Argument} & \textbf{Description} \\
\midrule
path & The file path to append to. \textit{(path-like)} \\
value ... & Optional explicit text values to append. When omitted, pipeline input becomes the appended text. (optional) \\
\bottomrule
\end{tabular}

\textbf{Pipeline input:} scalar, record
 --- When no explicit value arguments are supplied, pipeline values are rendered as plain text and appended to the file.

\textbf{Output:} Returns the resulting filesystem entry for the written file.

\begin{lstlisting}
append-file ./notes.txt " more"
echo "tail line" | append-file ./notes.txt
\end{lstlisting}

\begin{notebox}
These commands use ToSh's plain-text serialization rules, not the rich table renderer, so they are safe for intentional file output.
\end{notebox}

\end{cmdbox}

\subsection{\texttt{back}}
\label{ref:back}
\icmd{back}

\begin{cmdbox}{back}
Goes back to the previous directory in the stack.

\begin{signature}
back
\end{signature}

\textbf{Output:} The FileSystemEntry for the previous directory.

\begin{lstlisting}
back
\end{lstlisting}

\end{cmdbox}

\subsection{\texttt{basename}}
\label{ref:basename}
\icmd{basename}

\begin{cmdbox}{basename}
Returns the file name portion of each path.

\begin{signature}
basename <path> [suffix]
\end{signature}

\begin{tabular}{@{}>{\ttfamily\small\raggedright\arraybackslash}p{5cm} L{9cm}@{}}
\toprule
\textbf{Argument} & \textbf{Description} \\
\midrule
path & One or more paths to extract the filename from. \textit{(path-like)} \\
suffix & Optional suffix to strip from the filename. (optional) \\
\bottomrule
\end{tabular}

\textbf{Pipeline input:} scalar, list
 --- Accepts piped path-like values.

\textbf{Output:} Returns the filename component of each path.

\begin{lstlisting}
basename /home/user/file.txt
basename file.tar.gz .tar.gz  # Strip suffix
\end{lstlisting}

\end{cmdbox}

\subsection{\texttt{cat}}
\label{ref:cat}
\icmd{cat}

\begin{cmdbox}{cat}
Reads one or more files or piped text sources and emits their contents.

\begin{signature}
cat [-n|-b] [-s] [-E] [-T] [-A] [path ...|-]
\end{signature}

\begin{tabular}{@{}>{\ttfamily\small\raggedright\arraybackslash}p{5cm} L{9cm}@{}}
\toprule
\textbf{Argument} & \textbf{Description} \\
\midrule
path ...|\allowbreak - & One or more file paths to concatenate, or `-` to read piped text input explicitly. (optional) \textit{(path-like|string)} \\
\bottomrule
\end{tabular}

\begin{tabular}{@{}>{\ttfamily\small\raggedright\arraybackslash}p{5cm} L{9cm}@{}}
\toprule
\textbf{Flag / Option} & \textbf{Description} \\
\midrule
-n & Number every emitted line. \\
-b & Number only non-blank lines. \\
-s & Squeeze repeated blank lines into a single blank line. \\
\bottomrule
\end{tabular}

\textbf{Pipeline input:} scalar, record, list
 --- With no explicit paths, path-like pipeline values are treated as files when they all resolve to existing files; otherwise pipeline values are treated as text input. Use `-` explicitly when you want stdin-style text alongside file arguments.

\textbf{Output:} Returns plain text lines by default, or numbered record rows with `Number` and `Text` when `-n` or `-b` is used.

\begin{lstlisting}
cat README.md
echo alpha beta | cat -
cat -n README.md
\end{lstlisting}

\begin{notebox}
With no explicit paths, `cat` treats piped values as file paths only when every value resolves to an existing file. Otherwise it treats the pipeline as text input. Use `-` explicitly when mixing file paths with piped text.
\end{notebox}

\end{cmdbox}

\subsection{\texttt{cd}}
\label{ref:cd}
\icmd{cd}

\begin{cmdbox}{cd}
Changes the current directory for the Tosh session.

\begin{signature}
cd [path | - | +]
\end{signature}

\begin{tabular}{@{}>{\ttfamily\small\raggedright\arraybackslash}p{5cm} L{9cm}@{}}
\toprule
\textbf{Argument} & \textbf{Description} \\
\midrule
path & Absolute or relative directory path. Tilde expands to the home directory. (optional) \textit{(path-like)} \\
- & Go back to the previous directory (pop from stack). (optional) \\
+ & Go forward in the directory stack. (optional) \\
\bottomrule
\end{tabular}

\textbf{Output:} The FileSystemEntry for the new working directory.

\begin{lstlisting}
cd ~/projects  # Absolute path
cd ..  # Parent directory
cd -  # Previous directory
cd  # Home directory
\end{lstlisting}

\end{cmdbox}

\subsection{\texttt{chmod}}
\label{ref:chmod}
\icmd{chmod}

\begin{cmdbox}{chmod}
Changes file permission bits.

\begin{signature}
chmod [-R] <mode> <path> [path...]
\end{signature}

\begin{tabular}{@{}>{\ttfamily\small\raggedright\arraybackslash}p{5cm} L{9cm}@{}}
\toprule
\textbf{Argument} & \textbf{Description} \\
\midrule
mode & Permission mode string (e.g. 755, u+rw, a+x). \\
path & One or more files or directories. \textit{(path-like)} \\
\bottomrule
\end{tabular}

\begin{tabular}{@{}>{\ttfamily\small\raggedright\arraybackslash}p{5cm} L{9cm}@{}}
\toprule
\textbf{Flag / Option} & \textbf{Description} \\
\midrule
-R & Operate recursively on directories. \\
\bottomrule
\end{tabular}

\textbf{Output:} Returns FileSystemEntry objects with long display for each changed path.

\begin{lstlisting}
chmod +x script.sh
chmod 755 script.sh  # Octal mode
chmod u+rw,go-w file.txt  # Symbolic mode
chmod -R a+r ./docs  # Recursive
\end{lstlisting}

\end{cmdbox}

\subsection{\texttt{chown}}
\label{ref:chown}
\icmd{chown}

\begin{cmdbox}{chown}
Changes file owner and group.

\begin{signature}
chown [-R] <owner>[:group] <path> [path...]
\end{signature}

\begin{tabular}{@{}>{\ttfamily\small\raggedright\arraybackslash}p{5cm} L{9cm}@{}}
\toprule
\textbf{Argument} & \textbf{Description} \\
\midrule
owner[:group] & Owner and optional group specification. \\
path & One or more files or directories. \textit{(path-like)} \\
\bottomrule
\end{tabular}

\begin{tabular}{@{}>{\ttfamily\small\raggedright\arraybackslash}p{5cm} L{9cm}@{}}
\toprule
\textbf{Flag / Option} & \textbf{Description} \\
\midrule
-R & Operate recursively on directories. \\
\bottomrule
\end{tabular}

\textbf{Output:} Returns FileSystemEntry objects for each changed path.

\begin{lstlisting}
chown user:group file.txt
chown -R www-data /var/www  # Recursive change
\end{lstlisting}

\end{cmdbox}

\subsection{\texttt{close}}
\label{ref:close}
\icmd{close}

\begin{cmdbox}{close}
Closes one or more managed file or server handles.

\begin{signature}
close <handle> [handle...]
\end{signature}

\begin{tabular}{@{}>{\ttfamily\small\raggedright\arraybackslash}p{5cm} L{9cm}@{}}
\toprule
\textbf{Argument} & \textbf{Description} \\
\midrule
handle ... & One or more managed file handles to close. (optional) \\
\bottomrule
\end{tabular}

\textbf{Pipeline input:} record
 --- Consumes piped file handles when explicit handles are omitted.

\textbf{Output:} Closes the handles and does not emit pipeline output.

\begin{lstlisting}
close $handle
echo $handle | close
\end{lstlisting}

\begin{notebox}
These commands work with managed file handles returned by `open-file` or by `FileSystemEntry` methods like `OpenText()` and `OpenRead()`. `seek` returns the handle so you can keep piping through the stream workflow, while `copy-to` copies from one compatible handle into another.
\end{notebox}

\end{cmdbox}

\subsection{\texttt{copy-to}}
\label{ref:copy-to}
\icmd{copy-to}

\begin{cmdbox}{copy-to}
Copies the remaining contents of one managed file handle into another compatible handle.

\begin{signature}
copy-to [source] <target>
\end{signature}

\begin{tabular}{@{}>{\ttfamily\small\raggedright\arraybackslash}p{5cm} L{9cm}@{}}
\toprule
\textbf{Argument} & \textbf{Description} \\
\midrule
source & The readable managed file handle to copy from. (optional) \\
target & The writable managed file handle to copy into. \\
\bottomrule
\end{tabular}

\textbf{Pipeline input:} record
 --- Consumes a piped source handle when you only pass the target handle explicitly.

\textbf{Output:} Returns the number of bytes or text characters copied into the target handle.

\begin{lstlisting}
copy-to $source $target
$source | copy-to $target
\end{lstlisting}

\begin{notebox}
These commands work with managed file handles returned by `open-file` or by `FileSystemEntry` methods like `OpenText()` and `OpenRead()`. `seek` returns the handle so you can keep piping through the stream workflow, while `copy-to` copies from one compatible handle into another.
\end{notebox}

\end{cmdbox}

\subsection{\texttt{cp}}
\label{ref:cp}
\icmd{cp}

\begin{cmdbox}{cp}
Copies a file or directory.

\begin{signature}
cp [-r] [-f] [-n] [-p] [-u] <source> [source ...] <destination>
\end{signature}

\begin{tabular}{@{}>{\ttfamily\small\raggedright\arraybackslash}p{5cm} L{9cm}@{}}
\toprule
\textbf{Argument} & \textbf{Description} \\
\midrule
source & One or more source files or directories. \textit{(path-like)} \\
destination & The target path. \textit{(path-like)} \\
\bottomrule
\end{tabular}

\begin{tabular}{@{}>{\ttfamily\small\raggedright\arraybackslash}p{5cm} L{9cm}@{}}
\toprule
\textbf{Flag / Option} & \textbf{Description} \\
\midrule
-r & Copy directories recursively. \\
-f & Force overwrite of existing files. \\
-n & Do not overwrite existing files. \\
-p & Preserve file attributes such as timestamps. \\
-u & Only copy when the source is newer than the destination. \\
\bottomrule
\end{tabular}

\textbf{Output:} Returns FileInfo or DirectoryInfo objects for each copied target.

\begin{lstlisting}
cp file.txt backup.txt
cp -r src/ dst/  # Recursive directory copy
\end{lstlisting}

\end{cmdbox}

\subsection{\texttt{df}}
\label{ref:df}
\icmd{df}

\begin{cmdbox}{df}
Returns mounted file system usage information.

\textit{Aliases:} \code{mounts}

\begin{signature}
df [-hTl] [-t type[,type...]] [-x type[,type...]] [--total] [--output \\
\quad columns] [--show columns] [--hide columns] [--show-all] [path ...]
\end{signature}

\begin{tabular}{@{}>{\ttfamily\small\raggedright\arraybackslash}p{5cm} L{9cm}@{}}
\toprule
\textbf{Argument} & \textbf{Description} \\
\midrule
path ... & Optional paths used to resolve the containing mounted filesystem. (optional) \textit{(path-like)} \\
\bottomrule
\end{tabular}

\begin{tabular}{@{}>{\ttfamily\small\raggedright\arraybackslash}p{5cm} L{9cm}@{}}
\toprule
\textbf{Flag / Option} & \textbf{Description} \\
\midrule
-h & Accepts the familiar human-readable flag; ToSh sizes are already typed and human-friendly. \\
-T & Ensures the filesystem type column is visible. \\
-l & Restricts the output to local filesystems. \\
-t <type[,type...]> & Includes only matching filesystem types. \\
-x <type[,type...]> & Excludes matching filesystem types. \\
--total & Appends a typed aggregate total row. \\
--output <columns> & Selects which filesystem properties are rendered. \\
\bottomrule
\end{tabular}

\textbf{Pipeline input:} list
 --- Uses piped path-like values when explicit paths are omitted. Without path input, `df` lists the full mounted filesystem set.

\textbf{Output:} Produces typed filesystem usage objects, with optional aggregate totals.

\begin{lstlisting}
df --total
df -t ext4 --show FileSystem,Type,UsePercent,MountedOn
df . | get { FileSystem, MountedOn }
\end{lstlisting}

\end{cmdbox}

\subsection{\texttt{dirname}}
\label{ref:dirname}
\icmd{dirname}

\begin{cmdbox}{dirname}
Returns the directory portion of each path.

\begin{signature}
dirname <path> [path...]
\end{signature}

\begin{tabular}{@{}>{\ttfamily\small\raggedright\arraybackslash}p{5cm} L{9cm}@{}}
\toprule
\textbf{Argument} & \textbf{Description} \\
\midrule
path & One or more paths to extract the directory from. \textit{(path-like)} \\
\bottomrule
\end{tabular}

\textbf{Pipeline input:} scalar, list
 --- Accepts piped path-like values.

\textbf{Output:} Returns the directory component of each path.

\begin{lstlisting}
dirname /home/user/file.txt
\end{lstlisting}

\end{cmdbox}

\subsection{\texttt{dirs}}
\label{ref:dirs}
\icmd{dirs}

\begin{cmdbox}{dirs}
Shows or manages the directory stack.

\begin{signature}
dirs [goto <index> | remove <index> | clear]
\end{signature}

\begin{tabular}{@{}>{\ttfamily\small\raggedright\arraybackslash}p{5cm} L{9cm}@{}}
\toprule
\textbf{Argument} & \textbf{Description} \\
\midrule
subcommand & goto <index>, remove <index>, or clear. (optional) \\
\bottomrule
\end{tabular}

\textbf{Output:} Lists the current directory stack.

\begin{lstlisting}
dirs
dirs goto 2  # Navigate to stack entry
dirs clear  # Clear the directory stack
\end{lstlisting}

\end{cmdbox}

\subsection{\texttt{du}}
\label{ref:du}
\icmd{du}

\begin{cmdbox}{du}
Returns disk usage information for files and directories.

\textit{Aliases:} \code{disk-usage}, \code{usage}

\begin{signature}
du [-a] [-s] [-d depth] [-h] [-c] [-x] [--time] [--show columns] [--hide \\
\quad columns] [--show-all] [path ...]
\end{signature}

\begin{tabular}{@{}>{\ttfamily\small\raggedright\arraybackslash}p{5cm} L{9cm}@{}}
\toprule
\textbf{Argument} & \textbf{Description} \\
\midrule
path ... & Optional roots to measure. (optional) \textit{(path-like)} \\
\bottomrule
\end{tabular}

\begin{tabular}{@{}>{\ttfamily\small\raggedright\arraybackslash}p{5cm} L{9cm}@{}}
\toprule
\textbf{Flag / Option} & \textbf{Description} \\
\midrule
-a & Include file rows as well as directory summaries. \\
-s & Summarize each root instead of emitting recursive rows. \\
-d <depth> & Limit recursion depth. \\
-h & Accepts the familiar human-readable flag; ToSh sizes are already typed and human-friendly. \\
-c & Appends a typed grand total row. \\
-x & Stay on the same filesystem as each requested root. \\
--time & Include the latest modified timestamp for each emitted row. \\
\bottomrule
\end{tabular}

\textbf{Pipeline input:} list
 --- Uses piped path-like roots when explicit paths are omitted. Falls back to the current directory when neither are present.

\textbf{Output:} Produces typed path-usage objects with optional modified-time metadata and aggregate totals.

\begin{lstlisting}
du -s .
du -a -c --time
du -x ./projects | get { Name, Size, Modified }
\end{lstlisting}

\end{cmdbox}

\subsection{\texttt{exists}}
\label{ref:exists}
\icmd{exists}

\begin{cmdbox}{exists}
Checks whether each path exists.

\begin{signature}
exists <path> [path...]
\end{signature}

\end{cmdbox}

\subsection{\texttt{find}}
\label{ref:find}
\icmd{find}

\begin{cmdbox}{find}
Recursively finds file system entries.

\begin{signature}
find [path ...] [-name pattern] [-iname pattern] [-regex pattern] \\
\quad [-iregex pattern] [-type f|d|l] [-maxdepth n] [-mindepth n] [-size \\
\quad +/-size] [-mtime +/-days] [-newer-than duration] [-older-than \\
\quad duration] [-empty]
\end{signature}

\begin{tabular}{@{}>{\ttfamily\small\raggedright\arraybackslash}p{5cm} L{9cm}@{}}
\toprule
\textbf{Argument} & \textbf{Description} \\
\midrule
root ... & One or more filesystem roots to search. (optional) \textit{(path-like)} \\
\bottomrule
\end{tabular}

\begin{tabular}{@{}>{\ttfamily\small\raggedright\arraybackslash}p{5cm} L{9cm}@{}}
\toprule
\textbf{Flag / Option} & \textbf{Description} \\
\midrule
-name <pattern> & Filter by shell-style name pattern. \\
-regex <pattern> & Filter by .NET regex against the root-relative path. \\
-iregex <pattern> & Case-insensitive regex filter against the root-relative path. \\
-type <file|\allowbreak dir|\allowbreak link> & Filter by filesystem entry kind. \\
\bottomrule
\end{tabular}

\textbf{Pipeline input:} list
 --- Uses piped path-like roots when explicit roots are omitted. Falls back to the current directory when neither are present.

\textbf{Output:} Returns typed filesystem entries with rich metadata that flow naturally through the object pipeline.

\begin{lstlisting}
find . -name *.tosh
find . -regex ".*\\.tosh$"
find . -iregex ".*readme.*"
\end{lstlisting}

\end{cmdbox}

\subsection{\texttt{findmnt}}
\label{ref:findmnt}
\icmd{findmnt}

\begin{cmdbox}{findmnt}
Wraps the system findmnt utility, returning typed mounted-filesystem objects for JSON-backed invocations.

\begin{signature}
findmnt [-AflbR] [-S source] [-T target] [-M mountpoint] [-t types] [-O \\
\quad options] [-o columns] [path-or-device ...]
\end{signature}

\begin{tabular}{@{}>{\ttfamily\small\raggedright\arraybackslash}p{5cm} L{9cm}@{}}
\toprule
\textbf{Argument} & \textbf{Description} \\
\midrule
{}[path-or-device ...] & Optional mountpoints or source devices to match. (optional) \textit{(path-like|string)} \\
\bottomrule
\end{tabular}

\begin{tabular}{@{}>{\ttfamily\small\raggedright\arraybackslash}p{5cm} L{9cm}@{}}
\toprule
\textbf{Flag / Option} & \textbf{Description} \\
\midrule
-S <source> & Match a source device or source specification. \\
-T <target> & Match the filesystem that contains a target path. \\
-M <mountpoint> & Match a specific mountpoint. \\
-t <types> & Limit the result set to specific filesystem types. \\
-O <options> & Limit the result set to mounts with matching options. \\
-R & Include submounts for matching filesystems. \\
-U & Drop duplicate targets. \\
-l & Return a flattened list instead of a hierarchy. \\
-A & Disable built-in filters and include everything findmnt normally hides. \\
-b & Render size-oriented columns in raw bytes. \\
-D & Use a `df`-style display preset. \\
-I & Use a `df -i`-style inode display preset. \\
-o <columns> & Select findmnt-style output columns such as `TARGET,SOURCE,FSTYPE,OPTIONS`. \\
--output-all & Expose every selectable structured findmnt column. \\
\bottomrule
\end{tabular}

\textbf{Pipeline input:} none
 --- The structured `findmnt` builtin is explicit-arg-first and does not currently consume pipeline input.

\textbf{Output:} Returns reusable mounted-filesystem objects with nested child mounts when the underlying `findmnt --json` output is hierarchical.

\begin{lstlisting}
findmnt  # Browse the mounted-filesystem tree as typed objects
findmnt -l | where _.Target.StartsWith("/run")  # Flatten the mount tree and filter by target path
findmnt -o TARGET,SOURCE,FSTYPE  # Pick explicit findmnt-style display columns without changing the underlying objects
\end{lstlisting}

\begin{notebox}
ToSh wraps `findmnt --json --bytes --output-all` so mounted filesystems stay as typed objects in the pipeline while the default renderer can still show them as a tree with columns. The default result is hierarchical, so use `-l` when you want flat filtering in a pipeline, and shell-facing aliases like `FsType` and `FsRoot` match the visible column names. Output-format-only modes like `--pairs`, `--raw`, `--noheadings`, polling, and verification currently fall back to the external `findmnt` utility unchanged.
\end{notebox}

\end{cmdbox}

\subsection{\texttt{flush}}
\label{ref:flush}
\icmd{flush}

\begin{cmdbox}{flush}
Flushes one or more managed file handles.

\begin{signature}
flush <handle> [handle...]
\end{signature}

\begin{tabular}{@{}>{\ttfamily\small\raggedright\arraybackslash}p{5cm} L{9cm}@{}}
\toprule
\textbf{Argument} & \textbf{Description} \\
\midrule
handle ... & One or more managed file handles to flush. (optional) \\
\bottomrule
\end{tabular}

\textbf{Pipeline input:} record
 --- Consumes piped file handles when explicit handles are omitted.

\textbf{Output:} Flushes the handles and does not emit pipeline output.

\begin{lstlisting}
flush $handle
echo $handle | flush
\end{lstlisting}

\begin{notebox}
These commands work with managed file handles returned by `open-file` or by `FileSystemEntry` methods like `OpenText()` and `OpenRead()`. `seek` returns the handle so you can keep piping through the stream workflow, while `copy-to` copies from one compatible handle into another.
\end{notebox}

\end{cmdbox}

\subsection{\texttt{forward}}
\label{ref:forward}
\icmd{forward}

\begin{cmdbox}{forward}
Goes forward to the next directory in the stack.

\begin{signature}
forward
\end{signature}

\textbf{Output:} The FileSystemEntry for the next directory.

\begin{lstlisting}
forward
\end{lstlisting}

\end{cmdbox}

\subsection{\texttt{glob}}
\label{ref:glob}
\icmd{glob}

\begin{cmdbox}{glob}
Expands filesystem glob patterns.

\begin{signature}
glob [-a] <pattern> [pattern ...]
\end{signature}

\begin{tabular}{@{}>{\ttfamily\small\raggedright\arraybackslash}p{5cm} L{9cm}@{}}
\toprule
\textbf{Argument} & \textbf{Description} \\
\midrule
pattern & One or more glob patterns to expand. \\
\bottomrule
\end{tabular}

\begin{tabular}{@{}>{\ttfamily\small\raggedright\arraybackslash}p{5cm} L{9cm}@{}}
\toprule
\textbf{Flag / Option} & \textbf{Description} \\
\midrule
-a & Include hidden entries in results. \\
\bottomrule
\end{tabular}

\textbf{Pipeline input:} list
 --- Uses piped patterns when no arguments are given.

\textbf{Output:} Returns FileSystemEntry objects for each matched path.

\begin{lstlisting}
glob *.txt
glob -a **/*.cs  # Recursive with hidden files
\end{lstlisting}

\end{cmdbox}

\subsection{\texttt{is-dir}}
\label{ref:is-dir}
\icmd{is-dir}

\begin{cmdbox}{is-dir}
Checks whether each path is a directory.

\begin{signature}
is-dir <path> [path...]
\end{signature}

\end{cmdbox}

\subsection{\texttt{is-file}}
\label{ref:is-file}
\icmd{is-file}

\begin{cmdbox}{is-file}
Checks whether each path is a file.

\begin{signature}
is-file <path> [path...]
\end{signature}

\end{cmdbox}

\subsection{\texttt{is-link}}
\label{ref:is-link}
\icmd{is-link}

\begin{cmdbox}{is-link}
Checks whether each path is a symbolic link.

\begin{signature}
is-link <path> [path...]
\end{signature}

\end{cmdbox}

\subsection{\texttt{length}}
\label{ref:length}
\icmd{length}

\begin{cmdbox}{length}
Returns the current stream length for one or more managed file handles.

\begin{signature}
length [handle ...]
\end{signature}

\begin{tabular}{@{}>{\ttfamily\small\raggedright\arraybackslash}p{5cm} L{9cm}@{}}
\toprule
\textbf{Argument} & \textbf{Description} \\
\midrule
handle ... & One or more managed file handles whose current stream length should be reported. (optional) \\
\bottomrule
\end{tabular}

\textbf{Pipeline input:} record
 --- Consumes piped file handles when explicit handles are omitted.

\textbf{Output:} Returns the current underlying stream length for each supplied handle.

\begin{lstlisting}
length $handle
echo $handle | length
\end{lstlisting}

\begin{notebox}
These commands work with managed file handles returned by `open-file` or by `FileSystemEntry` methods like `OpenText()` and `OpenRead()`. `seek` returns the handle so you can keep piping through the stream workflow, while `copy-to` copies from one compatible handle into another.
\end{notebox}

\end{cmdbox}

\subsection{\texttt{ln}}
\label{ref:ln}
\icmd{ln}

\begin{cmdbox}{ln}
Creates hard links or symbolic links.

\begin{signature}
ln [-s] [-f] <target> <link-path>
\end{signature}

\begin{tabular}{@{}>{\ttfamily\small\raggedright\arraybackslash}p{5cm} L{9cm}@{}}
\toprule
\textbf{Argument} & \textbf{Description} \\
\midrule
target & The target path that the link will point to. \textit{(path-like)} \\
link-path & The path where the link will be created. \textit{(path-like)} \\
\bottomrule
\end{tabular}

\begin{tabular}{@{}>{\ttfamily\small\raggedright\arraybackslash}p{5cm} L{9cm}@{}}
\toprule
\textbf{Flag / Option} & \textbf{Description} \\
\midrule
-s & Create a symbolic link instead of a hard link. \\
-f & Remove existing destination files. \\
\bottomrule
\end{tabular}

\textbf{Output:} Returns a FileSystemEntry for the created link.

\begin{lstlisting}
ln original.txt hardlink.txt
ln -s /usr/bin/python3 ./python  # Create symbolic link
\end{lstlisting}

\end{cmdbox}

\subsection{\texttt{ls}}
\label{ref:ls}
\icmd{ls}

\begin{cmdbox}{ls}
Lists file system entries.

\begin{signature}
ls [-aAldRFihrSt] [-T [depth]] [--sort name|size|time] [--time \\
\quad modified|access|created] [--group-directories-first] [--tree [depth]] \\
\quad [--show columns] [--hide columns] [--show-all] [path ...]
\end{signature}

\begin{tabular}{@{}>{\ttfamily\small\raggedright\arraybackslash}p{5cm} L{9cm}@{}}
\toprule
\textbf{Argument} & \textbf{Description} \\
\midrule
path ... & Optional directories or files to list. (optional) \textit{(path-like)} \\
\bottomrule
\end{tabular}

\begin{tabular}{@{}>{\ttfamily\small\raggedright\arraybackslash}p{5cm} L{9cm}@{}}
\toprule
\textbf{Flag / Option} & \textbf{Description} \\
\midrule
-a & Include hidden entries. \\
-A & Include hidden entries while matching standard almost-all ls behavior. \\
-l & Use the long listing view. \\
-d & List directory arguments themselves instead of their contents. \\
-R & Traverse directories recursively. \\
-F & Classify names with shell-style suffixes like `*` and `@`. \\
-i & Include inode metadata in the compact table view. \\
-r & Reverse the current sort order. \\
-S & Sort by size descending. \\
-t & Sort by the active time field descending. \\
--sort <name|\allowbreak size|\allowbreak time> & Choose the primary listing sort field. \\
--time <modified|\allowbreak access|\allowbreak created> & Choose which time field long listings and time sorts use. \\
--group-directories-first & Group directories ahead of files before applying the primary sort. \\
-la & Combine hidden and long listing output. \\
\bottomrule
\end{tabular}

\textbf{Pipeline input:} none
 --- Ls is still explicit-arg-first; path input is not yet consumed from the pipeline.

\textbf{Output:} Produces typed filesystem entries that the display layer renders as shell tables by default.

\begin{lstlisting}
ls -la
ls -R --group-directories-first
ls -l --time access | where _.Type == file | get { Name, Accessed }
\end{lstlisting}

\begin{notebox}
Filesystem metadata stays typed in the pipeline, even when Tosh renders it like a shell table.
\end{notebox}

\end{cmdbox}

\subsection{\texttt{lsblk}}
\label{ref:lsblk}
\icmd{lsblk}

\begin{cmdbox}{lsblk}
Wraps the system lsblk utility, returning typed block-device objects for JSON-backed invocations.

\begin{signature}
lsblk [-aAdbflmpStzDNv] [-I list] [-e list] [-x column] [-o columns] \\
\quad [device ...]
\end{signature}

\begin{tabular}{@{}>{\ttfamily\small\raggedright\arraybackslash}p{5cm} L{9cm}@{}}
\toprule
\textbf{Argument} & \textbf{Description} \\
\midrule
{}[device ...] & Optional device paths or names to scope the query to. (optional) \textit{(path-like|string)} \\
\bottomrule
\end{tabular}

\begin{tabular}{@{}>{\ttfamily\small\raggedright\arraybackslash}p{5cm} L{9cm}@{}}
\toprule
\textbf{Flag / Option} & \textbf{Description} \\
\midrule
-a & Include empty devices. \\
-A & Hide empty devices. \\
-d & Suppress dependencies and child devices. \\
-b & Render size-oriented columns in raw bytes. \\
-f & Use the filesystem-oriented display preset. \\
-m & Use the permissions-oriented display preset. \\
-t & Use the topology-oriented display preset. \\
-D & Use the discard-oriented display preset. \\
-z & Use the zoned-device display preset. \\
-p & Show full device paths. \\
-S & Restrict the query to SCSI devices. \\
-N & Restrict the query to NVMe devices. \\
-v & Restrict the query to virtio devices. \\
-I <majors> & Include only the specified major numbers. \\
-e <majors> & Exclude the specified major numbers. \\
-x <column> & Sort by an lsblk column name such as `NAME`, `SIZE`, or `PATH`. \\
-o <columns> & Select lsblk-style output columns such as `NAME,SIZE,FSTYPE,MOUNTPOINTS`. \\
-O & Expose every selectable structured lsblk column. \\
\bottomrule
\end{tabular}

\textbf{Pipeline input:} none
 --- The structured `lsblk` builtin is explicit-arg-first and does not currently consume pipeline input.

\textbf{Output:} Returns reusable block-device objects with nested child devices when the underlying `lsblk --json` output is hierarchical.

\begin{lstlisting}
lsblk  # Browse block devices as a tree-with-columns object table
lsblk -l -f | where _.FsType == "ntfs"  # Flatten the block-device view and filter by filesystem type
lsblk -o NAME,PATH,SIZE  # Pick explicit lsblk-style display columns without changing the underlying objects
\end{lstlisting}

\begin{notebox}
ToSh wraps `lsblk --json --bytes --output-all` so block devices stay as typed objects in the pipeline while the default renderer can still show them as a tree with columns. The default result is hierarchical, so use `-l` when you want flat filtering in a pipeline, and shell-facing aliases like `FsType`, `FsVer`, and `FsAvail` match the visible column names. Output-format-only flags like `--pairs`, `--raw`, and `--noheadings` currently fall back to the external `lsblk` utility unchanged.
\end{notebox}

\end{cmdbox}

\subsection{\texttt{mkdir}}
\label{ref:mkdir}
\icmd{mkdir}

\begin{cmdbox}{mkdir}
Creates one or more directories.

\begin{signature}
mkdir [-p] <path> [path...]
\end{signature}

\begin{tabular}{@{}>{\ttfamily\small\raggedright\arraybackslash}p{5cm} L{9cm}@{}}
\toprule
\textbf{Argument} & \textbf{Description} \\
\midrule
path & One or more directory paths to create. \textit{(path-like)} \\
\bottomrule
\end{tabular}

\begin{tabular}{@{}>{\ttfamily\small\raggedright\arraybackslash}p{5cm} L{9cm}@{}}
\toprule
\textbf{Flag / Option} & \textbf{Description} \\
\midrule
-p & Create parent directories as needed; no error if existing. \\
\bottomrule
\end{tabular}

\textbf{Pipeline input:} list
 --- Accepts piped path-like values.

\textbf{Output:} Returns DirectoryInfo objects for each created directory.

\begin{lstlisting}
mkdir newdir
mkdir -p a/b/c  # Create nested directories
\end{lstlisting}

\end{cmdbox}

\subsection{\texttt{mkdir-temp}}
\label{ref:mkdir-temp}
\icmd{mkdir-temp}

\begin{cmdbox}{mkdir-temp}
Creates a temporary directory and returns it as a file system entry.

\begin{signature}
mkdir-temp [prefix]
\end{signature}

\begin{tabular}{@{}>{\ttfamily\small\raggedright\arraybackslash}p{5cm} L{9cm}@{}}
\toprule
\textbf{Argument} & \textbf{Description} \\
\midrule
prefix & Optional prefix for the directory name. Defaults to 'tosh'. (optional) \\
\bottomrule
\end{tabular}

\textbf{Output:} Returns a FileSystemEntry for the created temporary directory.

\begin{lstlisting}
mkdir-temp
mkdir-temp myapp  # Custom prefix
\end{lstlisting}

\end{cmdbox}

\subsection{\texttt{mv}}
\label{ref:mv}
\icmd{mv}

\begin{cmdbox}{mv}
Moves or renames files and directories.

\begin{signature}
mv [-nufTi] [-t directory] <source> [source ...] <destination>
\end{signature}

\begin{tabular}{@{}>{\ttfamily\small\raggedright\arraybackslash}p{5cm} L{9cm}@{}}
\toprule
\textbf{Argument} & \textbf{Description} \\
\midrule
source ... & One or more source paths to move. \textit{(path-like)} \\
destination & The destination path or directory. (optional) \textit{(path-like)} \\
\bottomrule
\end{tabular}

\begin{tabular}{@{}>{\ttfamily\small\raggedright\arraybackslash}p{5cm} L{9cm}@{}}
\toprule
\textbf{Flag / Option} & \textbf{Description} \\
\midrule
-n & Do not overwrite existing files. \\
-u & Only move when the source is newer than the destination. \\
-f & Force overwrite for file targets, clearing a previous `-n`. \\
-t <directory> & Use an explicit destination directory. \\
-T & Treat the destination as a normal path, not a target directory. \\
\bottomrule
\end{tabular}

\textbf{Pipeline input:} none
 --- The current `mv` implementation is explicit-arg-first and does not consume pipeline input.

\textbf{Output:} Returns FileInfo or DirectoryInfo objects for the moved targets.

\begin{lstlisting}
mv old.txt new.txt
mv -n *.txt archive/
mv -t archive alpha.txt beta.txt
\end{lstlisting}

\begin{notebox}
Mv now overwrites existing file targets by default, closer to Unix `mv`. `-n`, `-u`, `-t`, and `-T` are available to control that behavior explicitly.
\end{notebox}

\end{cmdbox}

\subsection{\texttt{open-file}}
\label{ref:open-file}
\icmd{open-file}

\begin{cmdbox}{open-file}
Opens one or more files as managed text or binary handles.

\begin{signature}
open-file [--read|--write|--append] [--binary] [--encoding <name>] \\
\quad <path> [path...]
\end{signature}

\begin{tabular}{@{}>{\ttfamily\small\raggedright\arraybackslash}p{5cm} L{9cm}@{}}
\toprule
\textbf{Argument} & \textbf{Description} \\
\midrule
path ... & One or more file paths to open. (optional) \textit{(path-like)} \\
\bottomrule
\end{tabular}

\begin{tabular}{@{}>{\ttfamily\small\raggedright\arraybackslash}p{5cm} L{9cm}@{}}
\toprule
\textbf{Flag / Option} & \textbf{Description} \\
\midrule
--read, -r & Open for reading. This is the default. \\
--write, -w & Open for writing and replace any previous contents. \\
--append, -a & Open for writing at the end of the file. \\
--binary, -b & Open a binary handle instead of a text handle. \\
--encoding <name> & Use a specific text encoding for text writers. \\
\bottomrule
\end{tabular}

\textbf{Pipeline input:} list
 --- Consumes piped path-like input when explicit file paths are omitted.

\textbf{Output:} Returns managed file-handle objects that can be passed to `read-from`, `read-line-from`, `read-to-end`, `write-to`, `write-line-to`, `flush`, `close`, `seek`, `position`, `length`, and `copy-to`.

\begin{lstlisting}
open-file ./notes.txt
open-file --write ./notes.txt
open-file --binary --append ./data.bin
\end{lstlisting}

\begin{notebox}
Open-file is the start of ToSh's managed stream system. It returns explicit text or binary handle objects instead of hiding file resources behind implicit properties. Active handles are also visible through `\$tosh.Session.OpenHandles` and `\$tosh.Session.OpenHandleCount`.
\end{notebox}

\end{cmdbox}

\subsection{\texttt{position}}
\label{ref:position}
\icmd{position}

\begin{cmdbox}{position}
Returns the current stream position for one or more managed file handles.

\begin{signature}
position [handle ...]
\end{signature}

\begin{tabular}{@{}>{\ttfamily\small\raggedright\arraybackslash}p{5cm} L{9cm}@{}}
\toprule
\textbf{Argument} & \textbf{Description} \\
\midrule
handle ... & One or more managed file handles whose current position should be reported. (optional) \\
\bottomrule
\end{tabular}

\textbf{Pipeline input:} record
 --- Consumes piped file handles when explicit handles are omitted.

\textbf{Output:} Returns the current byte position for each supplied handle when that position can be reported safely.

\begin{lstlisting}
position $handle
echo $handle | position
\end{lstlisting}

\begin{notebox}
These commands work with managed file handles returned by `open-file` or by `FileSystemEntry` methods like `OpenText()` and `OpenRead()`. `seek` returns the handle so you can keep piping through the stream workflow, while `copy-to` copies from one compatible handle into another.
\end{notebox}

\end{cmdbox}

\subsection{\texttt{pwd}}
\label{ref:pwd}
\icmd{pwd}

\begin{cmdbox}{pwd}
Returns the current directory as a DirectoryInfo object.

\begin{signature}
pwd
\end{signature}

\textbf{Output:} Returns the current working directory as a DirectoryInfo object.

\begin{lstlisting}
pwd
\end{lstlisting}

\end{cmdbox}

\subsection{\texttt{read-bytes}}
\label{ref:read-bytes}
\icmd{read-bytes}

\begin{cmdbox}{read-bytes}
Reads one or more files and returns each file as a byte array.

\begin{signature}
read-bytes <path> [path...]
\end{signature}

\begin{tabular}{@{}>{\ttfamily\small\raggedright\arraybackslash}p{5cm} L{9cm}@{}}
\toprule
\textbf{Argument} & \textbf{Description} \\
\midrule
path ... & One or more file paths to read as raw byte arrays. (optional) \textit{(path-like)} \\
\bottomrule
\end{tabular}

\textbf{Pipeline input:} list
 --- Consumes piped path-like input when explicit file paths are omitted.

\textbf{Output:} Returns one byte-array value per file.

\begin{lstlisting}
read-bytes ./image.bin
read-bytes ./image.bin | type-of
\end{lstlisting}

\begin{notebox}
These commands accept normal path-like values, including strings, FileInfo, and ToSh FileSystemEntry objects.
\end{notebox}

\end{cmdbox}

\subsection{\texttt{read-file}}
\label{ref:read-file}
\icmd{read-file}

\begin{cmdbox}{read-file}
Reads one or more files and returns each file as a single string value.

\begin{signature}
read-file <path> [path...]
\end{signature}

\begin{tabular}{@{}>{\ttfamily\small\raggedright\arraybackslash}p{5cm} L{9cm}@{}}
\toprule
\textbf{Argument} & \textbf{Description} \\
\midrule
path ... & One or more file paths to read as whole-text values. (optional) \textit{(path-like)} \\
\bottomrule
\end{tabular}

\textbf{Pipeline input:} list
 --- Consumes piped path-like input when explicit file paths are omitted.

\textbf{Output:} Returns one string value per file.

\begin{lstlisting}
read-file ./notes.txt
ls *.md | first | read-file
\end{lstlisting}

\begin{notebox}
These commands accept normal path-like values, including strings, FileInfo, and ToSh FileSystemEntry objects.
\end{notebox}

\end{cmdbox}

\subsection{\texttt{read-from}}
\label{ref:read-from}
\icmd{read-from}

\begin{cmdbox}{read-from}
Reads a text or binary chunk from an open managed file handle.

\begin{signature}
read-from [handle] [count]
\end{signature}

\begin{tabular}{@{}>{\ttfamily\small\raggedright\arraybackslash}p{5cm} L{9cm}@{}}
\toprule
\textbf{Argument} & \textbf{Description} \\
\midrule
handle & The managed file handle to read from. (optional) \\
count & Optional chunk size. Defaults to 4096. (optional) \textit{(int)} \\
\bottomrule
\end{tabular}

\textbf{Pipeline input:} record
 --- Consumes a piped file handle when no explicit handle argument is supplied.

\textbf{Output:} Returns a string chunk for text handles or a byte array chunk for binary handles.

\begin{lstlisting}
$handle | read-from 64
read-from $handle 4096
\end{lstlisting}

\begin{notebox}
These commands work with managed file handles returned by `open-file` or by `FileSystemEntry` methods like `OpenText()` and `OpenRead()`. `seek` returns the handle so you can keep piping through the stream workflow, while `copy-to` copies from one compatible handle into another.
\end{notebox}

\end{cmdbox}

\subsection{\texttt{read-line-from}}
\label{ref:read-line-from}
\icmd{read-line-from}

\begin{cmdbox}{read-line-from}
Reads the next text line from an open managed file handle.

\begin{signature}
read-line-from [handle]
\end{signature}

\begin{tabular}{@{}>{\ttfamily\small\raggedright\arraybackslash}p{5cm} L{9cm}@{}}
\toprule
\textbf{Argument} & \textbf{Description} \\
\midrule
handle & The managed text file handle to read from. (optional) \\
\bottomrule
\end{tabular}

\textbf{Pipeline input:} record
 --- Consumes a piped file handle when no explicit handle argument is supplied.

\textbf{Output:} Returns the next line as a ShellTextLine value, or nothing at end-of-file.

\begin{lstlisting}
$handle | read-line-from
read-line-from $handle
\end{lstlisting}

\begin{notebox}
These commands work with managed file handles returned by `open-file` or by `FileSystemEntry` methods like `OpenText()` and `OpenRead()`. `seek` returns the handle so you can keep piping through the stream workflow, while `copy-to` copies from one compatible handle into another.
\end{notebox}

\end{cmdbox}

\subsection{\texttt{read-lines}}
\label{ref:read-lines}
\icmd{read-lines}

\begin{cmdbox}{read-lines}
Reads one or more files and emits their contents line-by-line.

\begin{signature}
read-lines <path> [path...]
\end{signature}

\begin{tabular}{@{}>{\ttfamily\small\raggedright\arraybackslash}p{5cm} L{9cm}@{}}
\toprule
\textbf{Argument} & \textbf{Description} \\
\midrule
path ... & One or more file paths to read line-by-line. (optional) \textit{(path-like)} \\
\bottomrule
\end{tabular}

\textbf{Pipeline input:} list
 --- Consumes piped path-like input when explicit file paths are omitted.

\textbf{Output:} Returns ShellTextLine values, one per line across the supplied files.

\begin{lstlisting}
read-lines ./notes.txt
read-lines ./notes.txt | grep error
\end{lstlisting}

\begin{notebox}
These commands accept normal path-like values, including strings, FileInfo, and ToSh FileSystemEntry objects.
\end{notebox}

\end{cmdbox}

\subsection{\texttt{read-to-end}}
\label{ref:read-to-end}
\icmd{read-to-end}

\begin{cmdbox}{read-to-end}
Reads the remainder of an open managed file handle.

\begin{signature}
read-to-end [handle]
\end{signature}

\begin{tabular}{@{}>{\ttfamily\small\raggedright\arraybackslash}p{5cm} L{9cm}@{}}
\toprule
\textbf{Argument} & \textbf{Description} \\
\midrule
handle & The managed file handle to read from. (optional) \\
\bottomrule
\end{tabular}

\textbf{Pipeline input:} record
 --- Consumes a piped file handle when no explicit handle argument is supplied.

\textbf{Output:} Returns the remaining text or bytes from the handle.

\begin{lstlisting}
$handle | read-to-end
read-to-end $handle
\end{lstlisting}

\begin{notebox}
These commands work with managed file handles returned by `open-file` or by `FileSystemEntry` methods like `OpenText()` and `OpenRead()`. `seek` returns the handle so you can keep piping through the stream workflow, while `copy-to` copies from one compatible handle into another.
\end{notebox}

\end{cmdbox}

\subsection{\texttt{readlink}}
\label{ref:readlink}
\icmd{readlink}

\begin{cmdbox}{readlink}
Returns symbolic link targets.

\begin{signature}
readlink [-f] <path> [path...]
\end{signature}

\begin{tabular}{@{}>{\ttfamily\small\raggedright\arraybackslash}p{5cm} L{9cm}@{}}
\toprule
\textbf{Argument} & \textbf{Description} \\
\midrule
path & One or more symbolic link paths. \textit{(path-like)} \\
\bottomrule
\end{tabular}

\begin{tabular}{@{}>{\ttfamily\small\raggedright\arraybackslash}p{5cm} L{9cm}@{}}
\toprule
\textbf{Flag / Option} & \textbf{Description} \\
\midrule
-f & Resolve the final target, following chains of symbolic links. \\
\bottomrule
\end{tabular}

\textbf{Output:} Returns the target path of each symbolic link.

\begin{lstlisting}
readlink ./link
readlink -f ./chain  # Canonicalize
\end{lstlisting}

\end{cmdbox}

\subsection{\texttt{realpath}}
\label{ref:realpath}
\icmd{realpath}

\begin{cmdbox}{realpath}
Returns fully resolved absolute paths.

\begin{signature}
realpath <path> [path...]
\end{signature}

\begin{tabular}{@{}>{\ttfamily\small\raggedright\arraybackslash}p{5cm} L{9cm}@{}}
\toprule
\textbf{Argument} & \textbf{Description} \\
\midrule
path & One or more paths to resolve. \textit{(path-like)} \\
\bottomrule
\end{tabular}

\textbf{Output:} Returns fully resolved absolute path strings.

\begin{lstlisting}
realpath ./relative/path
\end{lstlisting}

\end{cmdbox}

\subsection{\texttt{rm}}
\label{ref:rm}
\icmd{rm}

\begin{cmdbox}{rm}
Removes files or directories.

\begin{signature}
rm [-r] [-f] [-i] <path> [path...]
\end{signature}

\begin{tabular}{@{}>{\ttfamily\small\raggedright\arraybackslash}p{5cm} L{9cm}@{}}
\toprule
\textbf{Argument} & \textbf{Description} \\
\midrule
path & One or more files or directories to remove. \textit{(path-like)} \\
\bottomrule
\end{tabular}

\begin{tabular}{@{}>{\ttfamily\small\raggedright\arraybackslash}p{5cm} L{9cm}@{}}
\toprule
\textbf{Flag / Option} & \textbf{Description} \\
\midrule
-r & Remove directories and their contents recursively. \\
-f & Ignore nonexistent files; never prompt. \\
-i & Prompt before each removal. \\
\bottomrule
\end{tabular}

\textbf{Pipeline input:} list
 --- Accepts piped path-like values.

\textbf{Output:} Returns a RemovalResult with total size and descendant tree.

\begin{lstlisting}
rm temp.txt
rm -rf build/  # Force recursive delete
\end{lstlisting}

\end{cmdbox}

\subsection{\texttt{seek}}
\label{ref:seek}
\icmd{seek}

\begin{cmdbox}{seek}
Moves an open managed file handle to a new position and returns the handle for continued piping.

\begin{signature}
seek [handle] <offset> [begin|current|end]
\end{signature}

\begin{tabular}{@{}>{\ttfamily\small\raggedright\arraybackslash}p{5cm} L{9cm}@{}}
\toprule
\textbf{Argument} & \textbf{Description} \\
\midrule
handle & The managed file handle to reposition. (optional) \\
offset & The byte offset to seek to or by. \textit{(long)} \\
origin & The seek origin: begin, current, or end. Defaults to begin. (optional) \\
\bottomrule
\end{tabular}

\textbf{Pipeline input:} record
 --- Consumes a piped file handle when no explicit handle argument is supplied.

\textbf{Output:} Moves the handle and returns it so you can continue piping into `read-from`, `read-line-from`, `read-to-end`, or `copy-to`.

\begin{lstlisting}
seek $handle 0 begin
$handle | seek 128 current
\end{lstlisting}

\begin{notebox}
These commands work with managed file handles returned by `open-file` or by `FileSystemEntry` methods like `OpenText()` and `OpenRead()`. `seek` returns the handle so you can keep piping through the stream workflow, while `copy-to` copies from one compatible handle into another.
\end{notebox}

\end{cmdbox}

\subsection{\texttt{stat}}
\label{ref:stat}
\icmd{stat}

\begin{cmdbox}{stat}
Returns detailed metadata for one or more paths.

\begin{signature}
stat [-L] [-f|--filesystem] [--show columns] [--hide columns] \\
\quad [--show-all] <path> [path...]
\end{signature}

\begin{tabular}{@{}>{\ttfamily\small\raggedright\arraybackslash}p{5cm} L{9cm}@{}}
\toprule
\textbf{Argument} & \textbf{Description} \\
\midrule
path & One or more paths to inspect. \textit{(path-like)} \\
\bottomrule
\end{tabular}

\begin{tabular}{@{}>{\ttfamily\small\raggedright\arraybackslash}p{5cm} L{9cm}@{}}
\toprule
\textbf{Flag / Option} & \textbf{Description} \\
\midrule
-L & Dereference symlinks before reading metadata. \\
-f & Return filesystem usage information for the containing mount instead of file-entry metadata. \\
--show <columns> & Select which properties are rendered. \\
--hide <columns> & Hide specific properties from the output. \\
--show-all & Display every available column. \\
\bottomrule
\end{tabular}

\textbf{Pipeline input:} list
 --- Uses piped path-like values when explicit paths are omitted.

\textbf{Output:} Typed filesystem-entry metadata, or filesystem-usage objects in -f mode.

\begin{lstlisting}
stat ./README.md
stat -L ./link.txt  # Dereference symlink
stat -f . | get { RequestedPath, FileSystem, MountedOn }  # Filesystem mode
\end{lstlisting}

\end{cmdbox}

\subsection{\texttt{tempfile}}
\label{ref:tempfile}
\icmd{tempfile}

\begin{cmdbox}{tempfile}
Creates a temporary file and returns it as a file system entry.

\begin{signature}
tempfile [prefix] [extension]
\end{signature}

\begin{tabular}{@{}>{\ttfamily\small\raggedright\arraybackslash}p{5cm} L{9cm}@{}}
\toprule
\textbf{Argument} & \textbf{Description} \\
\midrule
prefix & Optional prefix for the file name. Defaults to 'tosh'. (optional) \\
extension & Optional file extension. (optional) \\
\bottomrule
\end{tabular}

\textbf{Output:} Returns a FileSystemEntry for the created temporary file.

\begin{lstlisting}
tempfile
tempfile data .csv  # Custom prefix and extension
\end{lstlisting}

\end{cmdbox}

\subsection{\texttt{touch}}
\label{ref:touch}
\icmd{touch}

\begin{cmdbox}{touch}
Creates files or updates access and modification timestamps.

\begin{signature}
touch [-acm] [-d time|-r file] [-c] <path> [path...]
\end{signature}

\begin{tabular}{@{}>{\ttfamily\small\raggedright\arraybackslash}p{5cm} L{9cm}@{}}
\toprule
\textbf{Argument} & \textbf{Description} \\
\midrule
path ... & One or more filesystem paths to create or timestamp. \textit{(path-like)} \\
\bottomrule
\end{tabular}

\begin{tabular}{@{}>{\ttfamily\small\raggedright\arraybackslash}p{5cm} L{9cm}@{}}
\toprule
\textbf{Flag / Option} & \textbf{Description} \\
\midrule
-a & Update only the access time. \\
-m & Update only the modification time. \\
-c & Do not create missing files. \\
-d <time> & Use an explicit timestamp value. \\
-r <path> & Copy the timestamp from another file or directory. \\
\bottomrule
\end{tabular}

\textbf{Pipeline input:} list
 --- Consumes piped path-like input when explicit paths are omitted.

\textbf{Output:} Returns updated FileInfo or DirectoryInfo objects for the paths it touched.

\begin{lstlisting}
touch notes.txt
touch -c -a notes.txt
touch -d 2026-03-28T12:00:00 notes.txt
\end{lstlisting}

\begin{notebox}
Touch now supports `-a`, `-m`, `-c`, `-d`, and `-r`, plus grouped short flags like `-am`.
\end{notebox}

\end{cmdbox}

\subsection{\texttt{tree}}
\label{ref:tree}
\icmd{tree}

\begin{cmdbox}{tree}
Wraps the system tree utility, returning typed tree-entry objects.

\begin{signature}
tree [-adfL level] [--show columns] [--hide columns] [--show-all] [path \\
\quad ...]
\end{signature}

\begin{tabular}{@{}>{\ttfamily\small\raggedright\arraybackslash}p{5cm} L{9cm}@{}}
\toprule
\textbf{Argument} & \textbf{Description} \\
\midrule
path & Directory paths to tree. Defaults to current directory. (optional) \textit{(path-like)} \\
\bottomrule
\end{tabular}

\begin{tabular}{@{}>{\ttfamily\small\raggedright\arraybackslash}p{5cm} L{9cm}@{}}
\toprule
\textbf{Flag / Option} & \textbf{Description} \\
\midrule
-a & Include hidden files. \\
-d & List directories only. \\
-f & Print full path for each entry. \\
-L <level> & Limit recursion depth. \\
--show <columns> & Select which properties are rendered. \\
--hide <columns> & Hide specific properties from the output. \\
--show-all & Display every available column. \\
\bottomrule
\end{tabular}

\textbf{Output:} Returns typed tree-entry objects.

\begin{lstlisting}
tree
tree -L 2 -d  # Directories only, 2 levels
\end{lstlisting}

\end{cmdbox}

\subsection{\texttt{write-bytes}}
\label{ref:write-bytes}
\icmd{write-bytes}

\begin{cmdbox}{write-bytes}
Writes byte-oriented content to a file, replacing any previous contents.

\begin{signature}
write-bytes <path> [bytes...]
\end{signature}

\begin{tabular}{@{}>{\ttfamily\small\raggedright\arraybackslash}p{5cm} L{9cm}@{}}
\toprule
\textbf{Argument} & \textbf{Description} \\
\midrule
path & The file path to create or replace. \textit{(path-like)} \\
bytes ... & Optional explicit byte-oriented values. When omitted, pipeline input becomes the byte payload. (optional) \\
\bottomrule
\end{tabular}

\textbf{Pipeline input:} scalar, record
 --- When no explicit byte values are supplied, pipeline input is converted into a byte payload.

\textbf{Output:} Returns the resulting filesystem entry for the written file.

\begin{lstlisting}
write-bytes ./data.bin [1, 2, 3, 255]
read-bytes ./source.bin | write-bytes ./copy.bin
\end{lstlisting}

\begin{notebox}
Write-bytes accepts raw byte arrays, byte-like collections, and byte-convertible scalar values. Strings are encoded as UTF-8 in this first slice.
\end{notebox}

\end{cmdbox}

\subsection{\texttt{write-file}}
\label{ref:write-file}
\icmd{write-file}

\begin{cmdbox}{write-file}
Writes plain text to a file, replacing any previous contents.

\begin{signature}
write-file <path> [value...]
\end{signature}

\begin{tabular}{@{}>{\ttfamily\small\raggedright\arraybackslash}p{5cm} L{9cm}@{}}
\toprule
\textbf{Argument} & \textbf{Description} \\
\midrule
path & The file path to create or replace. \textit{(path-like)} \\
value ... & Optional explicit text values to write. When omitted, pipeline input becomes the file body. (optional) \\
\bottomrule
\end{tabular}

\textbf{Pipeline input:} scalar, record
 --- When no explicit value arguments are supplied, pipeline values are rendered as plain text and written to the file.

\textbf{Output:} Returns the resulting filesystem entry for the written file.

\begin{lstlisting}
write-file ./notes.txt "hello world"
echo alpha beta | write-file ./notes.txt
\end{lstlisting}

\begin{notebox}
These commands use ToSh's plain-text serialization rules, not the rich table renderer, so they are safe for intentional file output.
\end{notebox}

\end{cmdbox}

\subsection{\texttt{write-line-to}}
\label{ref:write-line-to}
\icmd{write-line-to}

\begin{cmdbox}{write-line-to}
Writes one or more text lines to an open managed text file handle.

\begin{signature}
write-line-to <handle> [value...]
\end{signature}

\begin{tabular}{@{}>{\ttfamily\small\raggedright\arraybackslash}p{5cm} L{9cm}@{}}
\toprule
\textbf{Argument} & \textbf{Description} \\
\midrule
handle & The managed text file handle to write into. \\
value ... & Optional explicit values to write as one line. When omitted, each pipeline value becomes its own line. (optional) \\
\bottomrule
\end{tabular}

\textbf{Pipeline input:} scalar, record
 --- When no explicit values are supplied, each pipeline value is written as its own line.

\textbf{Output:} Writes line-oriented text into the handle and does not emit pipeline output.

\begin{lstlisting}
write-line-to $handle hello world
echo alpha beta | write-line-to $handle
\end{lstlisting}

\begin{notebox}
These commands work with managed file handles returned by `open-file` or by `FileSystemEntry` methods like `OpenText()` and `OpenRead()`. `seek` returns the handle so you can keep piping through the stream workflow, while `copy-to` copies from one compatible handle into another.
\end{notebox}

\end{cmdbox}

\subsection{\texttt{write-to}}
\label{ref:write-to}
\icmd{write-to}

\begin{cmdbox}{write-to}
Writes plain text or bytes to an open managed file handle.

\begin{signature}
write-to <handle> [value...]
\end{signature}

\begin{tabular}{@{}>{\ttfamily\small\raggedright\arraybackslash}p{5cm} L{9cm}@{}}
\toprule
\textbf{Argument} & \textbf{Description} \\
\midrule
handle & The managed file handle to write into. \\
value ... & Optional explicit values to write. When omitted, pipeline input becomes the written payload. (optional) \\
\bottomrule
\end{tabular}

\textbf{Pipeline input:} scalar, record
 --- When no explicit values are supplied, pipeline values are written into the handle.

\textbf{Output:} Writes into the handle and does not emit pipeline output.

\begin{lstlisting}
write-to $handle hello world
echo alpha beta | write-to $handle
\end{lstlisting}

\begin{notebox}
These commands work with managed file handles returned by `open-file` or by `FileSystemEntry` methods like `OpenText()` and `OpenRead()`. `seek` returns the handle so you can keep piping through the stream workflow, while `copy-to` copies from one compatible handle into another.
\end{notebox}

\end{cmdbox}


\section{Functional}

\subsection{\texttt{compose}}
\label{ref:compose}
\icmd{compose}

\begin{cmdbox}{compose}
Composes two or more callables into a single callable that chains them left-to-right.

\begin{signature}
compose <callable1> <callable2> [callable ...]
\end{signature}

\end{cmdbox}

\subsection{\texttt{converge}}
\label{ref:converge}
\icmd{converge}

\begin{cmdbox}{converge}
Applies a callable to the seed repeatedly until two consecutive results are equal, then yields that stable value.

\begin{signature}
converge <seed> <callable|block>
\end{signature}

\end{cmdbox}

\subsection{\texttt{curry}}
\label{ref:curry}
\icmd{curry}

\begin{cmdbox}{curry}
Converts a fixed-arity callable into a curried callable value.

\begin{signature}
curry <callable>
\end{signature}

\begin{tabular}{@{}>{\ttfamily\small\raggedright\arraybackslash}p{5cm} L{9cm}@{}}
\toprule
\textbf{Argument} & \textbf{Description} \\
\midrule
callable & A fixed-arity callable value to curry. \\
\bottomrule
\end{tabular}

\textbf{Pipeline input:} none
 --- `curry` returns a callable that accumulates arguments until the original callable is saturated.

\textbf{Output:} Returns a curried callable. Calling it with too few arguments returns another callable; calling it with enough arguments runs the original callable.

\begin{lstlisting}
var add3 = func(a, b, c) => ($a + $b + $c); var curried = curry $add3  # Create a curried callable
var step1 = invoke $curried 1; var step2 = invoke $step1 2; invoke $step2 39  # Invoke a curried callable one step at a time
\end{lstlisting}

\end{cmdbox}

\subsection{\texttt{invoke}}
\label{ref:invoke}
\icmd{invoke}

\begin{cmdbox}{invoke}
Invokes a callable value such as a lambda or function object.

\begin{signature}
invoke <callable> [arg ...]
\end{signature}

\begin{tabular}{@{}>{\ttfamily\small\raggedright\arraybackslash}p{5cm} L{9cm}@{}}
\toprule
\textbf{Argument} & \textbf{Description} \\
\midrule
callable & A callable value such as an anonymous `func(...)` expression or other shell callable object. \\
arg ... & Optional positional arguments passed to the callable. (optional) \\
\bottomrule
\end{tabular}

\textbf{Pipeline input:} none
 --- `invoke` is explicit-argument based for now. It does not automatically feed the current pipeline into the callable.

\textbf{Output:} Returns whatever values the callable yields.

\begin{lstlisting}
invoke (func(x) => ($x * 2)) 21  # Invoke an inline anonymous function
var add = func(x, y) => ($x + $y); invoke $add 3 4  # Store a lambda and invoke it later
var factor = 3; var scale = func(x) => ($x * $factor); invoke $scale 7  # Invoke a closure that captures lexical state
\end{lstlisting}

\end{cmdbox}

\subsection{\texttt{iterate}}
\label{ref:iterate}
\icmd{iterate}

\begin{cmdbox}{iterate}
Generates an infinite sequence by repeatedly applying a callable to the previous result, starting from a seed. Pair with 'take' or 'take-while' to bound the output.

\begin{signature}
iterate <seed> <callable|block>
\end{signature}

\end{cmdbox}

\subsection{\texttt{partial}}
\label{ref:partial}
\icmd{partial}

\begin{cmdbox}{partial}
Binds leading arguments to a callable and returns a new callable value.

\begin{signature}
partial <callable> [arg ...]
\end{signature}

\begin{tabular}{@{}>{\ttfamily\small\raggedright\arraybackslash}p{5cm} L{9cm}@{}}
\toprule
\textbf{Argument} & \textbf{Description} \\
\midrule
callable & A callable value to partially apply. \\
arg ... & Leading positional arguments to bind ahead of time. (optional) \\
\bottomrule
\end{tabular}

\textbf{Pipeline input:} none
 --- `partial` is explicit-argument based and returns a new callable value.

\textbf{Output:} Returns a callable value that prepends the bound arguments before invoking the original callable.

\begin{lstlisting}
var add = func(x, y) => ($x + $y); var inc = partial $add 1; invoke $inc 41  # Bind the first argument of a callable
invoke (partial (func(a, b, c) => ($"{$a}-{$b}-{$c}")) alpha beta) gamma  # Bind multiple leading arguments
\end{lstlisting}

\end{cmdbox}

\subsection{\texttt{unfold}}
\label{ref:unfold}
\icmd{unfold}

\begin{cmdbox}{unfold}
Generates values from a seed by repeatedly applying a callable. The callable receives the current state and must return a [value, next-state] pair, or null to stop.

\begin{signature}
unfold <seed> <callable|block>
\end{signature}

\end{cmdbox}


\section{Network}

\subsection{\texttt{http}}
\label{ref:http}
\icmd{http}

\begin{cmdbox}{http}
Sends HTTP requests and returns structured response objects or decoded bodies.

\begin{signature}
http \\
\quad <get|\allowbreak post|\allowbreak put|\allowbreak patch|\allowbreak delete|\allowbreak head|\allowbreak options|\allowbreak request|\allowbreak send|\allowbreak serve|\allowbreak host|\allowbreak servers|\allowbreak stop> \\
\quad ...
\end{signature}

\begin{tabular}{@{}>{\ttfamily\small\raggedright\arraybackslash}p{5cm} L{9cm}@{}}
\toprule
\textbf{Argument} & \textbf{Description} \\
\midrule
<get|\allowbreak post|\allowbreak put|\allowbreak patch|\allowbreak delete|\allowbreak head|\allowbreak options> <url> & Send an HTTP request immediately. (optional) \\
request <method> <url> & Build an immutable HTTP request definition without sending it yet. (optional) \\
send [request] & Send an HttpRequestDefinition or HttpRequestMessage from an argument or the pipeline. (optional) \\
serve|\allowbreak host <dir> & Start a temporary HTTP file server rooted at a directory and return a live server handle. (optional) \\
servers & List open HTTP file server handles. (optional) \\
stop [handle|\allowbreak id ...] & Stop one or more HTTP file servers, or all of them when no target is provided. (optional) \\
\bottomrule
\end{tabular}

\begin{tabular}{@{}>{\ttfamily\small\raggedright\arraybackslash}p{5cm} L{9cm}@{}}
\toprule
\textbf{Flag / Option} & \textbf{Description} \\
\midrule
--header <name> <value> & Add a request header. Repeatable. \\
--json <value> & Serialize a value as JSON request content. \\
--body <text> & Send plain text request content. \\
--file <path> & Send request content from a file. \\
--form <record> & Send application/x-www-form-urlencoded content from a record-like value. \\
--content-type <value> & Override the request content type. \\
--timeout <duration> & Set the per-request timeout. \\
--bearer <token> & Add a Bearer authorization header. \\
--auth basic <user> <pass> & Add a Basic authorization header. \\
--follow |\allowbreak  --no-follow & Control redirect following for the request. \\
--as <response|\allowbreak json|\allowbreak text|\allowbreak bytes|\allowbreak lines> & Choose how the response should be materialized. \\
--out <path> & Write the raw response body bytes to a file. \\
--fail & Turn HTTP non-success status codes into diagnostics instead of returning a response object/body. \\
--browse & For `http serve`, render a lightweight browser page with directory listings and share metadata. \\
--upload & For `http serve`, accept uploads. Browser uploads and raw PUT/POST uploads are both supported. \\
--once & For `http serve`, close the server after the first handled request. \\
--bind <address> & Bind `http serve` to a specific address. \\
--lan & Bind `http serve` to all interfaces and advertise LAN-friendly share URLs for other devices. \\
--port <port> & Bind `http serve` to a specific port. Use `0` to request an ephemeral port. \\
--index <file> & Serve this index file for directories before directory listings. \\
--token <token> & Protect `http serve` with a fixed token. The returned ShareUrl includes it automatically. \\
--generate-token & Generate a random token for `http serve` and return it on the server handle. \\
\bottomrule
\end{tabular}

\textbf{Pipeline input:} scalar
 --- `http send` accepts a single HttpRequestDefinition or HttpRequestMessage from the pipeline.

\textbf{Output:} Returns either a structured HttpResponseInfo object or a decoded body, depending on `--as`. `http serve` returns live HttpFileServerHandle objects.

\begin{lstlisting}
http get https://example.com --as text  # Fetch a text response
http post https://example.com/api --json { Name = "Toast" } --as response  # Send JSON and keep the structured response
http request GET https://example.com | http send --as response  # Build then send a request object
http serve ./share --browse  # Start a temporary file server with a lightweight share page
http serve ./share --browse --lan  # Share a directory with other devices on the same network
http serve ./dropbox --upload --generate-token  # Start a temporary upload server with a generated access token
http servers | get { Id, ShareUrl, Upload }  # Inspect open temporary file servers
\end{lstlisting}

\begin{notebox}
The native `http` builtin is object-first and backed by .NET HttpClient. Use `--as response` for a structured response object, or `--as json|text|bytes|lines` to project the body directly. `http serve <dir>` starts a temporary file server and returns a live server handle; `--browse`, `--upload`, `--token`, `--generate-token`, and `--lan` turn it into a lightweight cross-platform sharing tool. Use `http servers`, `http stop`, or `close` to manage it.
\end{notebox}

\end{cmdbox}

\subsection{\texttt{ip}}
\label{ref:ip}
\icmd{ip}

\begin{cmdbox}{ip}
Wraps the system ip utility, returning typed objects for supported JSON-backed subcommands.

\begin{signature}
ip [addr|\allowbreak address|\allowbreak a|\allowbreak link|\allowbreak l|\allowbreak route|\allowbreak r [filter ...]] | ip <other-subcommand \\
\quad ...>
\end{signature}

\begin{tabular}{@{}>{\ttfamily\small\raggedright\arraybackslash}p{5cm} L{9cm}@{}}
\toprule
\textbf{Argument} & \textbf{Description} \\
\midrule
addr|\allowbreak address|\allowbreak a [filter ...] & Returns typed network-interface objects by invoking `ip -j addr` under the hood. (optional) \\
link|\allowbreak l [filter ...] & Returns typed network-interface link objects by invoking `ip -j link`. (optional) \\
route|\allowbreak r [filter ...] & Returns typed route objects by invoking `ip -j route`. (optional) \\
<other-subcommand ...> & For now, unsupported subcommands fall back to the system `ip` utility unchanged. (optional) \\
\bottomrule
\end{tabular}

\textbf{Pipeline input:} none
 --- The structured `ip` builtin is explicit-arg-first and does not currently consume pipeline input.

\textbf{Output:} For `ip addr` and `ip link`, returns typed interface objects with nested typed address objects where available. For `ip route`, returns typed route objects. Other subcommands currently pass through to the system `ip` utility's normal text output.

\begin{lstlisting}
ip addr  # List interfaces as structured objects
ip link | where _.IsUp  # Filter active interfaces in the pipeline
ip route | where _.IsDefault  # Show the default route as a typed object
ip addr | each { _.Addresses } | flatten | get Address  # Project nested typed IP addresses
\end{lstlisting}

\begin{notebox}
ToSh now wraps `ip addr`, `ip link`, and `ip route` around the JSON-capable system utility so the result flows through the pipeline as typed interface, address, and route objects. Other subcommands still fall back to the external `ip` utility unchanged.
\end{notebox}

\end{cmdbox}

\subsection{\texttt{ping}}
\label{ref:ping}
\icmd{ping}

\begin{cmdbox}{ping}
Pings a host and returns typed reply objects.

\begin{signature}
ping [-c count] [-W timeout-ms] <host>
\end{signature}

\end{cmdbox}


\section{Pipeline}

\subsection{\texttt{all}}
\label{ref:all}
\icmd{all}

\begin{cmdbox}{all}
Returns true if every pipeline value matches the predicate.

\begin{signature}
all <callable|block>
\end{signature}

\begin{tabular}{@{}>{\ttfamily\small\raggedright\arraybackslash}p{5cm} L{9cm}@{}}
\toprule
\textbf{Argument} & \textbf{Description} \\
\midrule
callable|\allowbreak block & A lambda or block predicate that returns boolean values. \\
\bottomrule
\end{tabular}

\textbf{Pipeline input:} scalar, record
 --- Consumes the current pipeline and stops at the first non-matching item.

\textbf{Output:} Returns `true` if every item matches the predicate; otherwise `false`.

\begin{lstlisting}
echo 2 4 6 | all func(x) => ((($x % 2) == 0))  # Check whether all items match
\end{lstlisting}

\end{cmdbox}

\subsection{\texttt{any}}
\label{ref:any}
\icmd{any}

\begin{cmdbox}{any}
Returns true if any pipeline value matches the predicate.

\begin{signature}
any <callable|block>
\end{signature}

\begin{tabular}{@{}>{\ttfamily\small\raggedright\arraybackslash}p{5cm} L{9cm}@{}}
\toprule
\textbf{Argument} & \textbf{Description} \\
\midrule
callable|\allowbreak block & A lambda or block predicate that returns boolean values. \\
\bottomrule
\end{tabular}

\textbf{Pipeline input:} scalar, record
 --- Consumes the current pipeline and stops at the first matching item.

\textbf{Output:} Returns `true` if any item matches the predicate; otherwise `false`.

\begin{lstlisting}
echo 1 2 3 | any func(x) => ($x == 2)  # Check whether any item matches
\end{lstlisting}

\end{cmdbox}

\subsection{\texttt{avg}}
\label{ref:avg}
\icmd{avg}

\begin{cmdbox}{avg}
Averages numeric, storage size, or timespan values.

\textit{Aliases:} \code{average}

\begin{signature}
avg [member-path]
\end{signature}

\end{cmdbox}

\subsection{\texttt{chunk}}
\label{ref:chunk}
\icmd{chunk}

\begin{cmdbox}{chunk}
Groups pipeline items into fixed-size batches.

\begin{signature}
chunk <size>
\end{signature}

\end{cmdbox}

\subsection{\texttt{collect}}
\label{ref:collect}
\icmd{collect}

\begin{cmdbox}{collect}
Collects all pipeline items into a single list.

\begin{signature}
collect
\end{signature}

\textbf{Pipeline input:} scalar, record, list, table
 --- Consumes the current pipeline and buffers every incoming item into one array result.

\textbf{Output:} Returns a single array containing the pipeline items in order.

\begin{lstlisting}
echo 1 2 3 | collect  # Collect scalar pipeline items into one array
ls *.cs | collect  # Capture a multi-item file listing as one value
findmnt -l | where _.FsType == ext4 | collect  # Buffer filtered structured rows into one array
\end{lstlisting}

\end{cmdbox}

\subsection{\texttt{count}}
\label{ref:count}
\icmd{count}

\begin{cmdbox}{count}
Counts the number of objects in the current pipeline.

\begin{signature}
count
\end{signature}

\end{cmdbox}

\subsection{\texttt{distinct}}
\label{ref:distinct}
\icmd{distinct}

\begin{cmdbox}{distinct}
Removes duplicate pipeline values.

\begin{signature}
distinct [member-path]
\end{signature}

\end{cmdbox}

\subsection{\texttt{each}}
\label{ref:each}
\icmd{each}

\begin{cmdbox}{each}
Executes a block or callable once for each input object.

\textit{Aliases:} \code{foreach}

\begin{signature}
each <callable|block>
\end{signature}

\begin{lstlisting}
echo one two | each { _.ToUpper() }
DriveInfo.GetDrives() | each func(d) => ($d.Name)
\end{lstlisting}

\begin{notebox}
Collections stay intact until you explicitly expand them with `each` or `flatten`.
\end{notebox}

\end{cmdbox}

\subsection{\texttt{filter}}
\label{ref:filter}
\icmd{filter}

\begin{cmdbox}{filter}
Filters pipeline values with a callable value or block predicate.

\begin{signature}
filter <callable|block>
\end{signature}

\begin{tabular}{@{}>{\ttfamily\small\raggedright\arraybackslash}p{5cm} L{9cm}@{}}
\toprule
\textbf{Argument} & \textbf{Description} \\
\midrule
callable|\allowbreak block & A lambda or block predicate that returns boolean values. \\
\bottomrule
\end{tabular}

\textbf{Pipeline input:} scalar, record
 --- Consumes the current pipeline and keeps only items whose predicate evaluates to true. Tree-shaped inputs are pruned like `where`.

\textbf{Output:} Returns the input items that matched the predicate.

\begin{lstlisting}
echo 1 2 3 4 | filter func(x) => ((($x % 2) == 0))  # Filter with a lambda
ls | filter { _.Type == file }  # Filter with a block
\end{lstlisting}

\end{cmdbox}

\subsection{\texttt{find-index}}
\label{ref:find-index}
\icmd{find-index}

\begin{cmdbox}{find-index}
Returns the 0-based index of the first pipeline value matching the predicate, or -1 if none match.

\begin{signature}
find-index <callable|block>
\end{signature}

\end{cmdbox}

\subsection{\texttt{first}}
\label{ref:first}
\icmd{first}

\begin{cmdbox}{first}
Returns the first object or first N objects from the pipeline.

\begin{signature}
first [count]
\end{signature}

\end{cmdbox}

\subsection{\texttt{flat-map}}
\label{ref:flat-map}
\icmd{flat-map}

\begin{cmdbox}{flat-map}
Transforms each pipeline value with a callable or block then flattens the results.

\begin{signature}
flat-map <callable|block>
\end{signature}

\end{cmdbox}

\subsection{\texttt{flatten}}
\label{ref:flatten}
\icmd{flatten}

\begin{cmdbox}{flatten}
Explicitly expands enumerable pipeline values by one level.

\begin{signature}
flatten
\end{signature}

\end{cmdbox}

\subsection{\texttt{frequencies}}
\label{ref:frequencies}
\icmd{frequencies}

\begin{cmdbox}{frequencies}
Counts occurrences of each distinct value in the pipeline.

\begin{signature}
frequencies
\end{signature}

\end{cmdbox}

\subsection{\texttt{get}}
\label{ref:get}
\icmd{get}

\begin{cmdbox}{get}
Gets a member, projects fields, or retrieves an item by index.

\begin{signature}
get <member-path> or get \{ <member-path>, ... \} or get <index>
\end{signature}

\begin{lstlisting}
ls -la | get Name
ps | get { Name, PID, Memory }
echo 1 2 3 | get func(x) => ($x * 2)
\end{lstlisting}

\end{cmdbox}

\subsection{\texttt{group-by}}
\label{ref:group-by}
\icmd{group-by}

\begin{cmdbox}{group-by}
Groups pipeline values by a member path, block, or callable.

\begin{signature}
group-by <member-path|callable|block>
\end{signature}

\begin{lstlisting}
ls | group-by Extension
ps | group-by func(p) => ($p.Name.Substring(0, 1))
\end{lstlisting}

\end{cmdbox}

\subsection{\texttt{group-while}}
\label{ref:group-while}
\icmd{group-while}

\begin{cmdbox}{group-while}
Groups consecutive items while the predicate holds.

\begin{signature}
group-while <callable|block>
\end{signature}

\end{cmdbox}

\subsection{\texttt{ignore}}
\label{ref:ignore}
\icmd{ignore}

\begin{cmdbox}{ignore}
Consumes and discards all pipeline input.

\begin{signature}
... | ignore
\end{signature}

\end{cmdbox}

\subsection{\texttt{inspect}}
\label{ref:inspect}
\icmd{inspect}

\begin{cmdbox}{inspect}
Inspects piped CLR objects inline, or returns legacy flat output with --flat.

\begin{signature}
inspect [-a] [--flat]
\end{signature}

\begin{lstlisting}
ls -la | first | inspect
new System.Random() | inspect -a
new System.Random() | inspect --flat
\end{lstlisting}

\begin{notebox}
Inspect opens an inline tree browser for CLR values in interactive sessions. Use `-a` for non-public/static members, `--flat` for the legacy static inspection object output, and `i` inside the browser to insert the selected member text into the active REPL line at the cursor (or queue it for the next prompt when no line is active). In the REPL, `F2` tries to inspect the reference under the cursor, and `Alt+I` is available as a fallback on terminals that do not expose function keys cleanly.
\end{notebox}

\end{cmdbox}

\subsection{\texttt{interleave}}
\label{ref:interleave}
\icmd{interleave}

\begin{cmdbox}{interleave}
Alternates items from the pipeline with items from another sequence.

\begin{signature}
interleave <other-sequence>
\end{signature}

\end{cmdbox}

\subsection{\texttt{last}}
\label{ref:last}
\icmd{last}

\begin{cmdbox}{last}
Returns the last object or last N objects from the pipeline.

\begin{signature}
last [count]
\end{signature}

\end{cmdbox}

\subsection{\texttt{map}}
\label{ref:map}
\icmd{map}

\begin{cmdbox}{map}
Transforms each pipeline value with a callable value or block.

\begin{signature}
map <callable|block>
\end{signature}

\begin{tabular}{@{}>{\ttfamily\small\raggedright\arraybackslash}p{5cm} L{9cm}@{}}
\toprule
\textbf{Argument} & \textbf{Description} \\
\midrule
callable|\allowbreak block & A lambda or block that transforms each input item into exactly one output value. \\
\bottomrule
\end{tabular}

\textbf{Pipeline input:} scalar, record
 --- Consumes the current pipeline and emits one transformed value per input item.

\textbf{Output:} Returns one transformed value for each input item.

\begin{lstlisting}
echo 1 2 3 | map func(x) => ($x * 2)  # Transform values with a lambda
ls | map { _.Name }  # Transform values with a block
\end{lstlisting}

\end{cmdbox}

\subsection{\texttt{max}}
\label{ref:max}
\icmd{max}

\begin{cmdbox}{max}
Returns the maximum pipeline value.

\begin{signature}
max [member-path]
\end{signature}

\end{cmdbox}

\subsection{\texttt{min}}
\label{ref:min}
\icmd{min}

\begin{cmdbox}{min}
Returns the minimum pipeline value.

\begin{signature}
min [member-path]
\end{signature}

\end{cmdbox}

\subsection{\texttt{none}}
\label{ref:none}
\icmd{none}

\begin{cmdbox}{none}
Returns true if no pipeline values match the predicate.

\begin{signature}
none <callable|block>
\end{signature}

\begin{tabular}{@{}>{\ttfamily\small\raggedright\arraybackslash}p{5cm} L{9cm}@{}}
\toprule
\textbf{Argument} & \textbf{Description} \\
\midrule
callable|\allowbreak block & A lambda or block predicate that returns boolean values. \\
\bottomrule
\end{tabular}

\textbf{Pipeline input:} scalar, record
 --- Consumes the current pipeline and stops at the first matching item.

\textbf{Output:} Returns `true` if no items match the predicate; otherwise `false`.

\begin{lstlisting}
echo 1 2 3 | none func(x) => ($x > 10)  # Check whether no items match
\end{lstlisting}

\end{cmdbox}

\subsection{\texttt{partition}}
\label{ref:partition}
\icmd{partition}

\begin{cmdbox}{partition}
Splits pipeline values into two lists based on a predicate: [matches, non-matches].

\begin{signature}
partition <callable|block>
\end{signature}

\end{cmdbox}

\subsection{\texttt{pick}}
\label{ref:pick}
\icmd{pick}

\begin{cmdbox}{pick}
Gets a member, projects fields, or retrieves an item by index.

\begin{signature}
pick <member-path> or pick \{ <member-path>, ... \} or pick <index>
\end{signature}

\end{cmdbox}

\subsection{\texttt{reduce}}
\label{ref:reduce}
\icmd{reduce}

\begin{cmdbox}{reduce}
Folds the current pipeline into one value using a seed and callable value or block.

\begin{signature}
reduce <seed> <callable|block>
\end{signature}

\begin{tabular}{@{}>{\ttfamily\small\raggedright\arraybackslash}p{5cm} L{9cm}@{}}
\toprule
\textbf{Argument} & \textbf{Description} \\
\midrule
seed & The initial accumulator value. \\
callable|\allowbreak block & A lambda or block that combines the current accumulator with each input item and returns the next accumulator. \\
\bottomrule
\end{tabular}

\textbf{Pipeline input:} scalar, record
 --- Consumes the current pipeline from left to right and folds it into one final value.

\textbf{Output:} Returns the final accumulator value. On an empty input stream, `reduce` returns the seed unchanged.

\begin{lstlisting}
echo 1 2 3 4 | reduce 0 func(acc, x) => ($acc + $x)  # Fold numeric values
echo one two three | reduce "" { $acc + _.Substring(0, 1) }  # Fold with a block
\end{lstlisting}

\end{cmdbox}

\subsection{\texttt{rename}}
\label{ref:rename}
\icmd{rename}

\begin{cmdbox}{rename}
Renames fields on record-like objects.

\begin{signature}
rename <old> <new> [old2 new2 ...]
\end{signature}

\end{cmdbox}

\subsection{\texttt{reverse}}
\label{ref:reverse}
\icmd{reverse}

\begin{cmdbox}{reverse}
Reverses the order of the current pipeline objects.

\begin{signature}
reverse
\end{signature}

\end{cmdbox}

\subsection{\texttt{scan}}
\label{ref:scan}
\icmd{scan}

\begin{cmdbox}{scan}
Like reduce but yields every intermediate accumulator value.

\begin{signature}
scan <seed> <callable|block>
\end{signature}

\end{cmdbox}

\subsection{\texttt{select}}
\label{ref:select}
\icmd{select}

\begin{cmdbox}{select}
Gets a member, projects fields, or retrieves an item by index.

\begin{signature}
select <member-path> or select \{ <member-path>, ... \} or select <index>
\end{signature}

\begin{lstlisting}
echo 1 2 3 | select func(x) => ($x * 2)
ls | select func(item) => ($item.Name)
\end{lstlisting}

\end{cmdbox}

\subsection{\texttt{skip}}
\label{ref:skip}
\icmd{skip}

\begin{cmdbox}{skip}
Skips the first object or first N objects from the pipeline.

\begin{signature}
skip [count]
\end{signature}

\end{cmdbox}

\subsection{\texttt{skip-until}}
\label{ref:skip-until}
\icmd{skip-until}

\begin{cmdbox}{skip-until}
Skips input values until the predicate becomes true.

\begin{signature}
skip-until <predicate-expression|callable>
\end{signature}

\end{cmdbox}

\subsection{\texttt{skip-while}}
\label{ref:skip-while}
\icmd{skip-while}

\begin{cmdbox}{skip-while}
Skips input values while the predicate remains true.

\begin{signature}
skip-while <predicate-expression|callable>
\end{signature}

\begin{lstlisting}
echo 1 2 3 4 | skip-while { _ < 3 }
echo 1 2 3 4 | skip-while func(x) => ($x < 3)
\end{lstlisting}

\end{cmdbox}

\subsection{\texttt{sort}}
\label{ref:sort}
\icmd{sort}

\begin{cmdbox}{sort}
Sorts the current pipeline objects.

\textit{Aliases:} \code{sort-by}

\begin{signature}
sort [-r|--reverse] [-n|--numeric] [-u|--unique] [-h|--human-numeric] \\
\quad [member-path|callable|block]
\end{signature}

\begin{tabular}{@{}>{\ttfamily\small\raggedright\arraybackslash}p{5cm} L{9cm}@{}}
\toprule
\textbf{Argument} & \textbf{Description} \\
\midrule
key & A member path, callable, or block to extract the sort key. (optional) \textit{(member-path|callable|block)} \\
\bottomrule
\end{tabular}

\begin{tabular}{@{}>{\ttfamily\small\raggedright\arraybackslash}p{5cm} L{9cm}@{}}
\toprule
\textbf{Flag / Option} & \textbf{Description} \\
\midrule
-r & Reverse the sort order. \\
-n & Use numeric comparison. \\
-u & Remove duplicate values after sorting. \\
-h & Human-numeric sort (understands storage sizes). \\
\bottomrule
\end{tabular}

\textbf{Pipeline input:} scalar
 --- Collects all pipeline objects, sorts them, then re-emits.

\textbf{Output:} The input pipeline objects in sorted order.

\begin{lstlisting}
ps | sort Memory  # Sort by property
ls -la | sort Modified | reverse  # Sort then reverse
ps | sort func(p) => ($p.Name.Length)  # Lambda sort key
\end{lstlisting}

\end{cmdbox}

\subsection{\texttt{sum}}
\label{ref:sum}
\icmd{sum}

\begin{cmdbox}{sum}
Sums numeric, storage size, or timespan values.

\begin{signature}
sum [member-path]
\end{signature}

\end{cmdbox}

\subsection{\texttt{summary}}
\label{ref:summary}
\icmd{summary}

\begin{cmdbox}{summary}
Computes structured summary objects for requested aggregates over the pipeline.

\textit{Aliases:} \code{summarize}

\begin{signature}
summary [--sum [columns]] [--avg [columns]] [--min [columns]] [--max \\
\quad [columns]] [--count [columns]]
\end{signature}

\begin{tabular}{@{}>{\ttfamily\small\raggedright\arraybackslash}p{5cm} L{9cm}@{}}
\toprule
\textbf{Argument} & \textbf{Description} \\
\midrule
{}[column|\allowbreak member-path] [--sum [columns]] [--avg [columns]] [--min [columns]] [--max [columns]] [--count [columns]] & Alias for `summarize`. (optional) \\
\bottomrule
\end{tabular}

\textbf{Pipeline input:} scalar, record, list, table
 --- Alias for `summarize`. Consumes the current pipeline rows and returns only ColumnSummary objects.

\textbf{Output:} Alias for `summarize`. Returns ColumnSummary objects only.

\begin{lstlisting}
summary _.Used
seq 5 | summary --count
ps | summary --avg Memory --max Memory
\end{lstlisting}

\end{cmdbox}

\subsection{\texttt{take-until}}
\label{ref:take-until}
\icmd{take-until}

\begin{cmdbox}{take-until}
Yields input values until the predicate becomes true.

\begin{signature}
take-until <predicate-expression|callable>
\end{signature}

\end{cmdbox}

\subsection{\texttt{take-while}}
\label{ref:take-while}
\icmd{take-while}

\begin{cmdbox}{take-while}
Yields input values while the predicate remains true.

\begin{signature}
take-while <predicate-expression|callable>
\end{signature}

\begin{lstlisting}
echo 1 2 3 4 | take-while { _ < 3 }
echo 1 2 3 4 | take-while func(x) => ($x < 3)
\end{lstlisting}

\end{cmdbox}

\subsection{\texttt{tee}}
\label{ref:tee}
\icmd{tee}

\begin{cmdbox}{tee}
Passes values through while also writing them out or capturing them.

\begin{signature}
tee [-a] [path] or tee -v <name>
\end{signature}

\end{cmdbox}

\subsection{\texttt{transpose}}
\label{ref:transpose}
\icmd{transpose}

\begin{cmdbox}{transpose}
Pivots rows into columns. Each input record's keys become column headers and values become rows.

\begin{signature}
transpose
\end{signature}

\end{cmdbox}

\subsection{\texttt{where}}
\label{ref:where}
\icmd{where}

\begin{cmdbox}{where}
Filters pipeline objects with a predicate block or callable.

\begin{signature}
where <predicate-expression|callable>
\end{signature}

\begin{tabular}{@{}>{\ttfamily\small\raggedright\arraybackslash}p{5cm} L{9cm}@{}}
\toprule
\textbf{Argument} & \textbf{Description} \\
\midrule
predicate & A predicate block or callable that returns a boolean. \textit{(block|callable)} \\
\bottomrule
\end{tabular}

\textbf{Pipeline input:} scalar
 --- Consumes any pipeline objects and tests each against the predicate.

\textbf{Output:} Pipeline objects for which the predicate returned true.

\begin{lstlisting}
ls -la | where _.Type == file  # Filter by property
ls -la | where func(item) => ($item.Name.ToLower().EndsWith(".md"))  # Lambda predicate
\end{lstlisting}

\begin{notebox}
Inside predicate expressions, bare member access resolves against the current pipeline object.
\end{notebox}

\end{cmdbox}

\subsection{\texttt{window}}
\label{ref:window}
\icmd{window}

\begin{cmdbox}{window}
Yields sliding windows of a given size over the pipeline.

\begin{signature}
window <size> [callable|block]
\end{signature}

\end{cmdbox}

\subsection{\texttt{xargs}}
\label{ref:xargs}
\icmd{xargs}

\begin{cmdbox}{xargs}
Builds command invocations from text input.

\begin{signature}
xargs [-n count] <command> [fixed-arg ...]
\end{signature}

\end{cmdbox}

\subsection{\texttt{zip}}
\label{ref:zip}
\icmd{zip}

\begin{cmdbox}{zip}
Merges two sequences pairwise. The second sequence is the first argument (an array).

\begin{signature}
zip <other-sequence> [callable|block]
\end{signature}

\end{cmdbox}


\section{Process}

\subsection{\texttt{jobs}}
\label{ref:jobs}
\icmd{jobs}

\begin{cmdbox}{jobs}
Lists ToSh background jobs.

\begin{signature}
jobs
\end{signature}

\begin{notebox}
Jobs lists ToSh background jobs started with a trailing `\&`. A background launch updates `\$tosh.Last.Result` with the started job info, while `jobs` and `wait-for` are the primary inspection commands.
\end{notebox}

\end{cmdbox}

\subsection{\texttt{kill}}
\label{ref:kill}
\icmd{kill}

\begin{cmdbox}{kill}
Stops a ToSh background job or operating-system process.

\begin{signature}
kill <job-id|pid> [job-id|pid ...]
\end{signature}

\begin{notebox}
Kill can stop either a ToSh background job or a native operating-system process by pid.
\end{notebox}

\end{cmdbox}

\subsection{\texttt{lsfd}}
\label{ref:lsfd}
\icmd{lsfd}

\begin{cmdbox}{lsfd}
Wraps the system lsfd utility, returning typed open-file-descriptor objects for JSON-backed invocations.

\begin{signature}
lsfd [-l] [-p pid[,pid...]] [-i[4|6]] [-o columns] \\
\quad [--summary[=only|\allowbreak append|\allowbreak never]]
\end{signature}

\begin{tabular}{@{}>{\ttfamily\small\raggedright\arraybackslash}p{5cm} L{9cm}@{}}
\toprule
\textbf{Argument} & \textbf{Description} \\
\midrule
{}[-p pid[,pid...]] & Optionally restrict the query to specific processes. (optional) \\
\bottomrule
\end{tabular}

\begin{tabular}{@{}>{\ttfamily\small\raggedright\arraybackslash}p{5cm} L{9cm}@{}}
\toprule
\textbf{Flag / Option} & \textbf{Description} \\
\midrule
-l & Include thread-level rows with a `TID` column. \\
-p <pid(s)> & Restrict the query to one or more process ids. \\
-i[4|\allowbreak 6] & Restrict the result set to IPv4 and/or IPv6 sockets. \\
-o <columns> & Select lsfd-style columns such as `COMMAND,PID,FD,TYPE,NAME`. \\
--summary[=only|\allowbreak append|\allowbreak never] & Include or isolate lsfd summary counters alongside structured descriptor rows. \\
--show <columns> & Use ToSh display-only column selection without changing the underlying `FileDescriptorInfo` objects. \\
--hide <columns> & Hide display columns while keeping the full typed rows in the pipeline. \\
--show-all & Expose every selectable structured lsfd column discoverable from the local `lsfd -H` catalog. \\
\bottomrule
\end{tabular}

\textbf{Pipeline input:} none
 --- The structured `lsfd` builtin is explicit-arg-first and does not currently consume pipeline input.

\textbf{Output:} Returns typed open-file-descriptor rows and, when summary mode is enabled, typed counter rows describing totals such as open files and sockets.

\begin{lstlisting}
lsfd  # Browse open file descriptors as typed rows
lsfd -p 1 -o COMMAND,PID,ASSOC,TYPE,NAME  # Inspect a specific process with explicit lsfd columns
lsfd --summary=only  # Show typed summary counters only
\end{lstlisting}

\begin{notebox}
ToSh wraps `lsfd --json` so open-file-descriptor rows stay typed in the pipeline. `--summary=only` yields typed counters, and `--summary=append` returns both row and summary objects. Text-format-only modes like `--raw`, `--noheadings`, filter expressions, and custom counters currently fall back to the external `lsfd` utility unchanged.
\end{notebox}

\end{cmdbox}

\subsection{\texttt{ps}}
\label{ref:ps}
\icmd{ps}

\begin{cmdbox}{ps}
Lists running processes as Tosh process objects.

\begin{signature}
ps [-eAf] [-p pid[,pid...]] [--ppid pid[,pid...]] [-u user[,user...]] \\
\quad [-t tty[,tty...]] [-o columns] [--sort field[,field...]] [--show \\
\quad columns] [--hide columns] [--show-all] [name-or-id ...]
\end{signature}

\begin{tabular}{@{}>{\ttfamily\small\raggedright\arraybackslash}p{5cm} L{9cm}@{}}
\toprule
\textbf{Argument} & \textbf{Description} \\
\midrule
name-or-id & Optional process names or PIDs to filter. (optional) \\
\bottomrule
\end{tabular}

\begin{tabular}{@{}>{\ttfamily\small\raggedright\arraybackslash}p{5cm} L{9cm}@{}}
\toprule
\textbf{Flag / Option} & \textbf{Description} \\
\midrule
-e & Show every process. \\
-A & Show every process (alias for -e). \\
-f & Full listing with extra columns. \\
-p <pid,...> & Select processes by PID. \\
--ppid <pid,...> & Select by parent PID. \\
-u <user,...> & Select by owning user. \\
-t <tty,...> & Select by terminal. \\
-o <columns> & Specify output columns. \\
--sort <field,...> & Sort by the given fields. \\
--show <columns> & Select which properties are rendered. \\
--hide <columns> & Hide specific properties from the output. \\
--show-all & Display every available column. \\
\bottomrule
\end{tabular}

\textbf{Output:} Typed process objects with properties like Name, Id, Memory, CPU, User.

\begin{lstlisting}
ps -f --sort -Id  # Full listing sorted by ID descending
ps -u root | get { Name, Id, User }  # Filter by user
ps | sort Memory | reverse | first 5  # Top 5 by memory
\end{lstlisting}

\begin{notebox}
Process memory values are surfaced as Tosh StorageSize objects, not raw strings.
\end{notebox}

\end{cmdbox}

\subsection{\texttt{signal}}
\label{ref:signal}
\icmd{signal}

\begin{cmdbox}{signal}
Sends a signal to a ToSh background job or operating-system process.

\begin{signature}
signal <signal> <job-id|pid> [job-id|pid ...]
\end{signature}

\begin{notebox}
Signal sends a named or numeric signal to a ToSh job or a native process id.
\end{notebox}

\end{cmdbox}

\subsection{\texttt{wait-for}}
\label{ref:wait-for}
\icmd{wait-for}

\begin{cmdbox}{wait-for}
Waits for one or more ToSh background jobs to finish.

\begin{signature}
wait-for [job-id ...]
\end{signature}

\begin{notebox}
Wait-for blocks until one or more background jobs finish and returns structured completion objects.
\end{notebox}

\end{cmdbox}


\section{Prompt}

\subsection{\texttt{prompt-dir}}
\label{ref:prompt-dir}
\icmd{prompt-dir}

\begin{cmdbox}{prompt-dir}
Returns the current directory as a styled prompt segment.

\begin{signature}
prompt-dir [--fg color] [--bg color] [--bold] [--depth n]
\end{signature}

\begin{lstlisting}
prompt-dir --fg blue --bold
prompt-dir --fg yellow --depth 2
\end{lstlisting}

\end{cmdbox}

\subsection{\texttt{prompt-duration}}
\label{ref:prompt-duration}
\icmd{prompt-duration}

\begin{cmdbox}{prompt-duration}
Returns the last command duration as a styled prompt segment.

\begin{signature}
prompt-duration [duration] [--fg color] [--bg color] [--bold] [--dim] \\
\quad [--threshold-ms value]
\end{signature}

\begin{lstlisting}
prompt-duration
prompt-duration 2.5s --threshold-ms 250
\end{lstlisting}

\end{cmdbox}

\subsection{\texttt{prompt-exit}}
\label{ref:prompt-exit}
\icmd{prompt-exit}

\begin{cmdbox}{prompt-exit}
Returns the last non-zero exit code as a styled prompt segment.

\begin{signature}
prompt-exit [code] [--fg color] [--bg color] [--bold] [--dim]
\end{signature}

\begin{lstlisting}
prompt-exit
prompt-exit 7 --fg red --bold
\end{lstlisting}

\end{cmdbox}

\subsection{\texttt{prompt-git}}
\label{ref:prompt-git}
\icmd{prompt-git}

\begin{cmdbox}{prompt-git}
Returns git branch and status as styled prompt segments.

\begin{signature}
prompt-git [--fg color] [--bg color] [--bold]
\end{signature}

\begin{lstlisting}
prompt-git
prompt-git --fg bright-green --bold
\end{lstlisting}

\end{cmdbox}

\subsection{\texttt{prompt-history}}
\label{ref:prompt-history}
\icmd{prompt-history}

\begin{cmdbox}{prompt-history}
Returns the next history id as a styled prompt segment.

\begin{signature}
prompt-history [id] [--fg color] [--bg color] [--bold] [--dim]
\end{signature}

\begin{lstlisting}
prompt-history
prompt-history 432 --fg gray --dim
\end{lstlisting}

\end{cmdbox}

\subsection{\texttt{prompt-jobs}}
\label{ref:prompt-jobs}
\icmd{prompt-jobs}

\begin{cmdbox}{prompt-jobs}
Returns the current background job count as a styled prompt segment.

\begin{signature}
prompt-jobs [count] [--fg color] [--bg color] [--bold] [--dim]
\end{signature}

\begin{lstlisting}
prompt-jobs
prompt-jobs 3 --fg yellow --bold
\end{lstlisting}

\end{cmdbox}

\subsection{\texttt{prompt-newline}}
\label{ref:prompt-newline}
\icmd{prompt-newline}

\begin{cmdbox}{prompt-newline}
Inserts a line break in the prompt.

\begin{signature}
prompt-newline
\end{signature}

\begin{lstlisting}
prompt-newline
\end{lstlisting}

\end{cmdbox}

\subsection{\texttt{prompt-text}}
\label{ref:prompt-text}
\icmd{prompt-text}

\begin{cmdbox}{prompt-text}
Returns literal text as a styled prompt segment.

\begin{signature}
prompt-text <text> [--fg color] [--bg color] [--bold] [--dim]
\end{signature}

\begin{lstlisting}
prompt-text "> " --fg cyan
prompt-text "::" --fg gray --dim
\end{lstlisting}

\end{cmdbox}

\subsection{\texttt{prompt-time}}
\label{ref:prompt-time}
\icmd{prompt-time}

\begin{cmdbox}{prompt-time}
Returns the current time as a styled prompt segment.

\begin{signature}
prompt-time [--fg color] [--bg color] [--bold] [--dim] [--format \\
\quad pattern]
\end{signature}

\begin{lstlisting}
prompt-time --dim
prompt-time --format "HH:mm:ss" --fg gray
\end{lstlisting}

\end{cmdbox}

\subsection{\texttt{prompt-userhost}}
\label{ref:prompt-userhost}
\icmd{prompt-userhost}

\begin{cmdbox}{prompt-userhost}
Returns the current user and host as a styled prompt segment.

\begin{signature}
prompt-userhost [--fg color] [--bg color] [--bold] [--dim]
\end{signature}

\begin{lstlisting}
prompt-userhost --dim
prompt-userhost --fg gray
\end{lstlisting}

\end{cmdbox}

\subsection{\texttt{styled}}
\label{ref:styled}
\icmd{styled}

\begin{cmdbox}{styled}
Creates a styled text segment with color and formatting.

\begin{signature}
styled <text> [--fg color] [--bg color] [--bold] [--italic] \\
\quad [--underline] [--dim]
\end{signature}

\begin{lstlisting}
styled "hello" --fg cyan --bold
styled "warning" --fg yellow --bg red
\end{lstlisting}

\end{cmdbox}


\section{Scripting}

\subsection{\texttt{assert}}
\label{ref:assert}
\icmd{assert}

\begin{cmdbox}{assert}
Asserts that a condition is true; throws a diagnostic error if it is false.

\begin{signature}
assert <predicate> [message]
\end{signature}

\end{cmdbox}

\subsection{\texttt{raise}}
\label{ref:raise}
\icmd{raise}

\begin{cmdbox}{raise}
Raises an event, invoking all registered handlers.

\begin{signature}
raise <event> | <event> | raise
\end{signature}

\end{cmdbox}


\section{Shell}

\subsection{\texttt{apropos}}
\label{ref:apropos}
\icmd{apropos}

\begin{cmdbox}{apropos}
Searches Tosh help topics with fuzzy matching.

\begin{signature}
apropos <query>
\end{signature}

\begin{lstlisting}
apropos json
apropos loop
\end{lstlisting}

\begin{notebox}
Apropos performs fuzzy help search across commands and Tosh language topics.
\end{notebox}

\end{cmdbox}

\subsection{\texttt{clear}}
\label{ref:clear}
\icmd{clear}

\begin{cmdbox}{clear}
Clears the terminal display.

\begin{signature}
clear
\end{signature}

\end{cmdbox}

\subsection{\texttt{config}}
\label{ref:config}
\icmd{config}

\begin{cmdbox}{config}
Gets or changes shell configuration.

\begin{signature}
config [browse [query]|get <path>|set <path> [value]|reset [path]|init \\
\quad [path]|reload]
\end{signature}

\begin{tabular}{@{}>{\ttfamily\small\raggedright\arraybackslash}p{5cm} L{9cm}@{}}
\toprule
\textbf{Argument} & \textbf{Description} \\
\midrule
browse [query] & Opens the full-screen config browser, optionally filtered by an initial query, with staged editing, subtree diffs, structured section and collection editors, reusable confirmation and validation surfaces, filesystem browsing, apply/save flows, startup reload/init actions, live prompt/style/theme previews, and raw text editing for advanced cases. (optional) \\
get <path> & Reads one config value. (optional) \\
set <path> [value] & Sets one config value. (optional) \\
reset [section] & Resets a section or the whole config object. (optional) \\
reload & Replays startup config files into the current session. (optional) \\
init [directory] & Scaffolds a new config directory. (optional) \\
\bottomrule
\end{tabular}

\textbf{Pipeline input:} scalar
 --- Only `config set` consumes piped scalar input, using it as the new value when no explicit value argument is present.

\textbf{Output:} Produces the live config object, one config value, status rows, or an interactive browser request depending on the form.

\end{cmdbox}

\subsection{\texttt{events}}
\label{ref:events}
\icmd{events}

\begin{cmdbox}{events}
Lists, inspects, and manages event handlers.

\begin{signature}
events [list|names|handlers <event>|remove <event> <handler>|clear \\
\quad <event>]
\end{signature}

\end{cmdbox}

\subsection{\texttt{exec}}
\label{ref:exec}
\icmd{exec}

\begin{cmdbox}{exec}
Replace the current ToSh process with an external command.

\begin{signature}
exec [--] <command> [arg ...]
\end{signature}

\begin{lstlisting}
exec tosh
exec zsh
exec /bin/sh -c "echo hi"
\end{lstlisting}

\begin{notebox}
Exec replaces the current ToSh process with an external command. On Unix-like systems it uses native process replacement, so `exec tosh` or `exec zsh` behaves like the shell built-in you may know from zsh.
\end{notebox}

\end{cmdbox}

\subsection{\texttt{exit}}
\label{ref:exit}
\icmd{exit}

\begin{cmdbox}{exit}
Requests the current Tosh session to exit.

\begin{signature}
exit
\end{signature}

\end{cmdbox}

\subsection{\texttt{help}}
\label{ref:help}
\icmd{help}

\begin{cmdbox}{help}
Shows searchable Tosh help for commands, language topics, CLR types, and externals.

\begin{signature}
help [--cli] [topic ... | browse [query] | search <query> | related \\
\quad <topic> | categories]
\end{signature}

\begin{tabular}{@{}>{\ttfamily\small\raggedright\arraybackslash}p{5cm} L{9cm}@{}}
\toprule
\textbf{Argument} & \textbf{Description} \\
\midrule
topic & The command, language topic, type, or external executable to describe. (optional) \\
--cli [query] & Opens the inline fuzzy tree browser, optionally seeded with an initial query or topic. (optional) \\
browse [query] & Opens the full-screen help browser, optionally filtered by an initial query. (optional) \\
search <query> & Searches help topics by name, alias, category, and description. (optional) \\
related <topic> & Shows related topics for the given command or language feature. (optional) \\
categories & Lists help categories and topic counts. (optional) \\
\bottomrule
\end{tabular}

\textbf{Pipeline input:} scalar
 --- With no explicit args, piped scalar values are treated as help topics. `search` and `related` also consume piped query/topic text.

\textbf{Output:} Produces help summaries, full help topics, search results, category rows, launches the inline browser with `--cli`, or returns an interactive browser request for `help browse` depending on the form.

\begin{lstlisting}
help ls  # Describe a command
help search regex  # Search help
help --cli regex  # Open the inline fuzzy help browser
help browse grep  # Open the help browser
\end{lstlisting}

\begin{notebox}
Use `help search <query>` or `apropos <query>` to find commands and language topics quickly, `help --cli` for the inline fuzzy tree browser, or `help browse` for the fullscreen split-pane browser. In the REPL, `F1` opens the inline help browser seeded from the token under the cursor, and `Alt+H` is available as a fallback on terminals that do not expose function keys cleanly.
\end{notebox}

\end{cmdbox}

\subsection{\texttt{history}}
\label{ref:history}
\icmd{history}

\begin{cmdbox}{history}
Shows, searches, expands, runs, deletes, saves, reloads, or clears shell history.

\begin{signature}
history [status|path|search <text>|expand <spec>|run <spec>|delete \\
\quad <spec>|save|reload|clear]
\end{signature}

\begin{tabular}{@{}>{\ttfamily\small\raggedright\arraybackslash}p{5cm} L{9cm}@{}}
\toprule
\textbf{Argument} & \textbf{Description} \\
\midrule
search <text> & Searches history entries by text. (optional) \\
expand <spec> & Expands a history event specification without running it. (optional) \\
run <spec> & Resolves and executes a history event specification. (optional) \\
delete <spec> & Deletes one or more history entries by id or spec. (optional) \\
path|\allowbreak save|\allowbreak reload|\allowbreak clear & History maintenance subcommands. (optional) \\
\bottomrule
\end{tabular}

\textbf{Pipeline input:} none
 --- History is producer-oriented; replay and expansion are explicit subcommands, while `!` syntax remains REPL-only sugar.

\textbf{Output:} Produces structured history entries, file paths, expanded command text, or replay results depending on the chosen subcommand.

\begin{notebox}
History is file-backed in normal interactive sessions and each entry now has a stable id. Use `history path`, `history search <text>`, `history delete <spec>`, `history save`, `history reload`, `history clear`, `history expand 237`, or `history run 237` to inspect or replay it. In the REPL, `Ctrl+R`, `!!`, `!237`, `!-2`, `!prefix`, `!?text?`, `!\$`, `!\^{}`, `!*`, and `\^{}old\^{}new\^{}` also work as interactive history features.
\end{notebox}

\end{cmdbox}

\subsection{\texttt{history-search}}
\label{ref:history-search}
\icmd{history-search}

\begin{cmdbox}{history-search}
Searches shell history entries by text.

\begin{signature}
history-search <text>
\end{signature}

\end{cmdbox}

\subsection{\texttt{read-line}}
\label{ref:read-line}
\icmd{read-line}

\begin{cmdbox}{read-line}
Reads a line of text from standard input.

\begin{signature}
read-line [prompt]
\end{signature}

\end{cmdbox}

\subsection{\texttt{tui}}
\label{ref:tui}
\icmd{tui}

\begin{cmdbox}{tui}
Interactive TUI components for scripts. Provides list pickers, confirmations, text input, file pickers, and custom screens. Use --cli for inline (non-fullscreen) prompts.

\begin{signature}
tui \\
\quad pick|\allowbreak confirm|\allowbreak input|\allowbreak file|\allowbreak filter|\allowbreak screen|\allowbreak add-list|\allowbreak add-text|\allowbreak add-input|\allowbreak add-picker|\allowbreak layout|\allowbreak run \\
\quad [options]
\end{signature}

\end{cmdbox}

\subsection{\texttt{view}}
\label{ref:view}
\icmd{view}

\begin{cmdbox}{view}
Gets or sets shell display preferences.

\begin{signature}
view \\
\quad [compact|\allowbreak detail|\allowbreak datetime|\allowbreak datetimeoffset|\allowbreak dateonly|\allowbreak timeonly|\allowbreak timespan|\allowbreak size|\allowbreak permissions|\allowbreak attributes|\allowbreak columns]
\end{signature}

\begin{lstlisting}
view dateonly scalar relative
view timeonly table 24h
view duration table seconds
\end{lstlisting}

\end{cmdbox}

\subsection{\texttt{which}}
\label{ref:which}
\icmd{which}

\begin{cmdbox}{which}
Resolves built-in commands and external executables.

\textit{Aliases:} \code{whence}

\begin{signature}
which <name ...>
\end{signature}

\end{cmdbox}


\section{System}

\subsection{\texttt{date}}
\label{ref:date}
\icmd{date}

\begin{cmdbox}{date}
Creates, parses, and adjusts date/time values.

\begin{signature}
date [-d|--date-only] [-t|--time-only] \\
\quad <now|\allowbreak utc-now|\allowbreak today|\allowbreak tomorrow|\allowbreak yesterday|\allowbreak parse|\allowbreak from-unix|\allowbreak from-unix-ms|\allowbreak date-only|\allowbreak time-only|\allowbreak <iso-date>> \\
\quad ... or <date> | date [-d] [-t] <add|\allowbreak sub|\allowbreak date-only|\allowbreak time-only> ...
\end{signature}

\begin{lstlisting}
date now -d -t
date parse 2026-03-29T12:34:56Z -d
date parse 2026-03-29T12:34:56Z | cast timeonly
\end{lstlisting}

\end{cmdbox}

\subsection{\texttt{env}}
\label{ref:env}
\icmd{env}

\begin{cmdbox}{env}
Lists, queries, sets, or unsets environment variables. Runs nested commands with temporary environment changes when a command follows.

\begin{signature}
env [name ...] | env [-u name] [name=value ...] | env [-u name] \\
\quad [name=value ...] -- <command ...>
\end{signature}

\begin{tabular}{@{}>{\ttfamily\small\raggedright\arraybackslash}p{5cm} L{9cm}@{}}
\toprule
\textbf{Argument} & \textbf{Description} \\
\midrule
name ... & With no assignments, query one or more environment variable names. (optional) \\
name=value ... & Temporary environment assignments for `env` output or a nested command. (optional) \\
command ... & Optional nested command to run under the temporary environment snapshot. (optional) \\
\bottomrule
\end{tabular}

\begin{tabular}{@{}>{\ttfamily\small\raggedright\arraybackslash}p{5cm} L{9cm}@{}}
\toprule
\textbf{Flag / Option} & \textbf{Description} \\
\midrule
-u <name> & Unset a variable in the temporary environment snapshot. \\
-- & Treat the remaining arguments as the nested command even when there are no assignments. \\
\bottomrule
\end{tabular}

\textbf{Pipeline input:} scalar, record
 --- With no explicit names, piped scalar values are treated as queried environment-variable names.

\textbf{Output:} Returns typed environment-variable entries, or the nested command's output when assignments/unsets are used with a command. Use `\$env.NAME` for direct value-only access when you want the variable value instead of the structured `env` entry object. Missing `\$env` members resolve to `null` rather than raising a missing-member error.

\begin{lstlisting}
env PATH  # Inspect the structured PATH entry
echo $env.PATH  # Read just the PATH value
echo $env.EDITOR  # Read another environment variable directly
\end{lstlisting}

\begin{notebox}
Env keeps its object-returning query mode, but it can also build a temporary environment snapshot with `name=value` and `-u name`, optionally running a nested command under that snapshot.
\end{notebox}

\end{cmdbox}

\subsection{\texttt{export}}
\label{ref:export}
\icmd{export}

\begin{cmdbox}{export}
Exports a Tosh value to the process environment.

\begin{signature}
export <name> [value]
\end{signature}

\end{cmdbox}

\subsection{\texttt{free}}
\label{ref:free}
\icmd{free}

\begin{cmdbox}{free}
Returns system memory and swap usage.

\begin{signature}
free
\end{signature}

\end{cmdbox}

\subsection{\texttt{guid}}
\label{ref:guid}
\icmd{guid}

\begin{cmdbox}{guid}
Creates, parses, formats, and inspects GUID values.

\begin{signature}
guid [new [v4|v7]|empty|parse|format <d|n|b|p|x>|info] [value ...]
\end{signature}

\begin{lstlisting}
guid
guid new v7
guid info 550e8400-e29b-41d4-a716-446655440000
\end{lstlisting}

\end{cmdbox}

\subsection{\texttt{hostname}}
\label{ref:hostname}
\icmd{hostname}

\begin{cmdbox}{hostname}
Returns the current host name.

\begin{signature}
hostname
\end{signature}

\end{cmdbox}

\subsection{\texttt{hostnamectl}}
\label{ref:hostnamectl}
\icmd{hostnamectl}

\begin{cmdbox}{hostnamectl}
Wraps hostnamectl, returning a structured host-status object for supported JSON-backed status queries.

\begin{signature}
hostnamectl [status | <other-command ...>]
\end{signature}

\begin{tabular}{@{}>{\ttfamily\small\raggedright\arraybackslash}p{5cm} L{9cm}@{}}
\toprule
\textbf{Argument} & \textbf{Description} \\
\midrule
{}[status] & With no subcommand, ToSh treats `hostnamectl` as a structured status query. Explicit `status` behaves the same way. (optional) \\
<other-command ...> & Mutating or unsupported commands fall back to the native `hostnamectl` utility unchanged. (optional) \\
\bottomrule
\end{tabular}

\begin{tabular}{@{}>{\ttfamily\small\raggedright\arraybackslash}p{5cm} L{9cm}@{}}
\toprule
\textbf{Flag / Option} & \textbf{Description} \\
\midrule
--show <columns> & Use ToSh display-only column selection on the structured host-status object. \\
--hide <columns> & Hide display columns while preserving the underlying typed host object. \\
--show-all & Expose every selectable structured host-status display column. \\
\bottomrule
\end{tabular}

\textbf{Pipeline input:} none
 --- The structured `hostnamectl` builtin is explicit-arg-first and does not currently consume pipeline input.

\textbf{Output:} Returns a structured host-status object for supported JSON-backed `hostnamectl status` queries. Other commands currently fall back to the native `hostnamectl` output.

\begin{lstlisting}
hostnamectl  # Inspect structured host status
hostnamectl --show Hostname,OperatingSystem,KernelRelease  # Render a focused host summary
hostnamectl | get { Hostname, MachineID, BootID }  # Project host identity properties in the pipeline
\end{lstlisting}

\end{cmdbox}

\subsection{\texttt{id}}
\label{ref:id}
\icmd{id}

\begin{cmdbox}{id}
Returns current user and group identity information.

\begin{signature}
id [-u|-g|-G] [-n]
\end{signature}

\end{cmdbox}

\subsection{\texttt{journalctl}}
\label{ref:journalctl}
\icmd{journalctl}

\begin{cmdbox}{journalctl}
Wraps journalctl, returning typed journal-entry objects for supported JSON-backed query invocations.

\begin{signature}
journalctl [-n count] [-u unit] [--since when] [--until when] [query \\
\quad ...]
\end{signature}

\begin{tabular}{@{}>{\ttfamily\small\raggedright\arraybackslash}p{5cm} L{9cm}@{}}
\toprule
\textbf{Argument} & \textbf{Description} \\
\midrule
{}[query ...] & Structured journal queries accept common `journalctl` filters such as match expressions, unit filters, priorities, boot selectors, and time windows. (optional) \\
\bottomrule
\end{tabular}

\begin{tabular}{@{}>{\ttfamily\small\raggedright\arraybackslash}p{5cm} L{9cm}@{}}
\toprule
\textbf{Flag / Option} & \textbf{Description} \\
\midrule
-n <count>|\allowbreak --lines <count> & Limit the structured result set to the most recent entries. \\
-u <unit>|\allowbreak --unit <unit> & Restrict structured output to a specific systemd unit. \\
--since <when> & Restrict entries to those at or after the supplied time. \\
--until <when> & Restrict entries to those at or before the supplied time. \\
-p <range>|\allowbreak --priority <range> & Restrict entries to a priority or priority range. \\
-b [id]|\allowbreak --boot [id] & Restrict entries to the current boot or a selected boot. \\
-g <pattern>|\allowbreak --grep <pattern> & Restrict structured output to entries whose messages match a pattern. \\
--user & Query the user journal instead of the system journal. \\
--system & Query the system journal explicitly. \\
--show <columns> & Use ToSh display-only column selection on the structured journal-entry rows. \\
--hide <columns> & Hide display columns while preserving the underlying structured journal entries. \\
--show-all & Expose every selectable structured journal-entry display column. \\
\bottomrule
\end{tabular}

\textbf{Pipeline input:} none
 --- The structured `journalctl` builtin is explicit-arg-first and does not currently consume pipeline input.

\textbf{Output:} Streams structured journal-entry objects from `journalctl -o json` for supported non-follow query paths. Unsupported or text-oriented modes fall back to the native utility.

\begin{lstlisting}
journalctl -n 20  # Stream the newest journal entries as structured objects
journalctl -u sshd.service | first 10  # Inspect a unit's logs in the pipeline
journalctl --since yesterday | where _.Priority <= 3  # Filter recent high-priority entries
\end{lstlisting}

\end{cmdbox}

\subsection{\texttt{loginctl}}
\label{ref:loginctl}
\icmd{loginctl}

\begin{cmdbox}{loginctl}
Wraps loginctl, returning typed session/user/seat rows and structured login property sets for supported queries.

\begin{signature}
loginctl [list-sessions | list-users | list-seats | show-session <id \\
\quad ...> | show-user <user ...> | show-seat <seat ...> | <other-subcommand \\
\quad ...>]
\end{signature}

\begin{tabular}{@{}>{\ttfamily\small\raggedright\arraybackslash}p{5cm} L{9cm}@{}}
\toprule
\textbf{Argument} & \textbf{Description} \\
\midrule
{}[list-sessions|\allowbreak list-users|\allowbreak list-seats] & With no subcommand, ToSh treats `loginctl` as a structured `list-sessions` query. The explicit list subcommands return typed rows. (optional) \\
show-session <id ...>|\allowbreak show-user <user ...>|\allowbreak show-seat <seat ...> & Returns structured login property sets for supported `show-*` queries. (optional) \\
<other-subcommand ...> & Unsupported subcommands fall back to the native `loginctl` utility unchanged. (optional) \\
\bottomrule
\end{tabular}

\begin{tabular}{@{}>{\ttfamily\small\raggedright\arraybackslash}p{5cm} L{9cm}@{}}
\toprule
\textbf{Flag / Option} & \textbf{Description} \\
\midrule
-p <property[,property...]>|\allowbreak --property <property[,property...]> & Restrict `show-*` to specific fetched properties. ToSh still injects the relevant identity property internally so multiple-result output stays structured. \\
--all & Include empty properties in supported `show-*` queries when the underlying `loginctl` invocation supports it. \\
--show <columns> & Use ToSh display-only column selection on structured list rows or property sets. \\
--hide <columns> & Hide display columns while preserving the underlying typed objects. \\
--show-all & Expose every selectable structured display column for the current output shape. \\
\bottomrule
\end{tabular}

\textbf{Pipeline input:} none
 --- The structured `loginctl` builtin is explicit-arg-first and does not currently consume pipeline input.

\textbf{Output:} Returns typed session, user, or seat rows for supported list queries, and structured property-set objects for supported `show-*` queries. Other subcommands currently fall back to the native `loginctl` output.

\begin{lstlisting}
loginctl  # List sessions as typed rows
loginctl list-users | where _.State == active  # Filter active login users in the pipeline
loginctl show-user 1000 | get { UID, Name, State, Sessions }  # Inspect structured user-login properties
\end{lstlisting}

\end{cmdbox}

\subsection{\texttt{lscpu}}
\label{ref:lscpu}
\icmd{lscpu}

\begin{cmdbox}{lscpu}
Wraps the system lscpu utility, returning typed CPU summary, topology, and cache objects for JSON-backed invocations.

\begin{signature}
lscpu [-B] [--hierarchic when] | lscpu -e[-all|-online|-offline] [-x] \\
\quad [-y] [-o columns] | lscpu -C [-B] [-o columns]
\end{signature}

\begin{tabular}{@{}>{\ttfamily\small\raggedright\arraybackslash}p{5cm} L{9cm}@{}}
\toprule
\textbf{Argument} & \textbf{Description} \\
\midrule
{}[-e|\allowbreak -C] & Without a mode flag, `lscpu` returns a structured CPU summary. `-e` switches to per-CPU topology rows and `-C` switches to CPU cache rows. (optional) \\
\bottomrule
\end{tabular}

\begin{tabular}{@{}>{\ttfamily\small\raggedright\arraybackslash}p{5cm} L{9cm}@{}}
\toprule
\textbf{Flag / Option} & \textbf{Description} \\
\midrule
-B & Render cache sizes in raw bytes for summary and cache views. \\
-e & Return per-CPU topology rows from `lscpu --extended --json`. \\
-C & Return CPU cache rows from `lscpu --caches --json`. \\
-a & Include both online and offline CPUs in extended mode. \\
-b & Restrict extended mode to online CPUs. \\
-c & Restrict extended mode to offline CPUs. \\
-x & Request hexadecimal CPU masks where the underlying `lscpu` mode supports them. \\
-y & Request physical IDs instead of logical IDs in extended mode. \\
-o <columns> & Select lscpu-style columns such as `CPU,NODE,SOCKET,CORE,ONLINE,MHZ` for `-e` or `NAME,LEVEL,TYPE,ONE-SIZE` for `-C`. \\
--output-all & Expose every selectable structured column for `-e` or `-C`. \\
--hierarchic <when> & Pass through `auto`, `always`, or `never` for the summary view. \\
\bottomrule
\end{tabular}

\textbf{Pipeline input:} none
 --- The structured `lscpu` builtin is explicit-arg-first and does not currently consume pipeline input.

\textbf{Output:} Returns a structured CPU summary by default, typed per-CPU topology rows with `-e`, or typed CPU cache rows with `-C`.

\begin{lstlisting}
lscpu  # Show the structured CPU summary
lscpu -e | where _.Online == true | first 8  # Browse structured per-CPU topology rows
lscpu -C -B | get { Name, Level, OneSize, AllSize }  # Inspect cache metadata with byte-oriented sizes
\end{lstlisting}

\begin{notebox}
ToSh wraps the JSON-capable `lscpu` modes instead of scraping text output. The default command yields a structured CPU summary, `-e` yields per-CPU topology rows, and `-C` yields cache rows. `--parse`, raw-only, help, version, and sysroot modes currently fall back to the external `lscpu` utility unchanged.
\end{notebox}

\end{cmdbox}

\subsection{\texttt{lsipc}}
\label{ref:lsipc}
\icmd{lsipc}

\begin{cmdbox}{lsipc}
Wraps the system lsipc utility, returning structured IPC resource and limit records for JSON-backed invocations.

\begin{signature}
lsipc [-m|-M|-q|-Q|-s|-S] [-g] [-i id|-N name] [-c] [-t] [-b] [-P] [-o \\
\quad columns]
\end{signature}

\begin{tabular}{@{}>{\ttfamily\small\raggedright\arraybackslash}p{5cm} L{9cm}@{}}
\toprule
\textbf{Argument} & \textbf{Description} \\
\midrule
{}[-m|\allowbreak -M|\allowbreak -q|\allowbreak -Q|\allowbreak -s|\allowbreak -S] & Select a specific IPC resource family. With no resource flag, `lsipc` returns the global IPC limits and usage summary. (optional) \\
\bottomrule
\end{tabular}

\begin{tabular}{@{}>{\ttfamily\small\raggedright\arraybackslash}p{5cm} L{9cm}@{}}
\toprule
\textbf{Flag / Option} & \textbf{Description} \\
\midrule
-m & Return System V shared-memory rows. \\
-M & Return POSIX shared-memory rows. \\
-q & Return System V message-queue rows. \\
-Q & Return POSIX message-queue rows. \\
-s & Return System V semaphore rows. \\
-S & Return POSIX semaphore rows. \\
-g & Return global usage/limit rows, optionally scoped by `-m`, `-q`, or `-s`. \\
-i <id> & Restrict the query to a specific System V IPC id. \\
-N <name> & Restrict the query to a specific POSIX IPC name. \\
-c & Include creator-related fields such as creator uid, user, and group. \\
-t & Include time-oriented fields such as attach, detach, change, or last-operation timestamps. \\
-b & Request byte-oriented numeric sizes from the underlying `lsipc` command. \\
-P & Render permissions numerically instead of symbolically. \\
-l & Force list output shape where the underlying `lsipc` mode supports it. \\
-o <columns> & Select lsipc-style columns such as `KEY,ID,OWNER,SIZE,NATTCH` or `RESOURCE,LIMIT,USED,USE\%`. \\
--show <columns> & Use ToSh display-only column selection on the structured rows after parsing. \\
--hide <columns> & Hide display columns while keeping the full structured rows in the pipeline. \\
\bottomrule
\end{tabular}

\textbf{Pipeline input:} none
 --- The structured `lsipc` builtin is explicit-arg-first and does not currently consume pipeline input.

\textbf{Output:} Returns structured IPC resource rows or global IPC limit/usage rows, with typed sizes, counts, and ISO-normalized timestamps where the underlying data supports them.

\begin{lstlisting}
lsipc  # Browse global IPC limits and current usage as structured rows
lsipc -m | first 5  # Inspect System V shared-memory rows
lsipc -g -m | get { Resource, Limit, Used, UsePercent }  # Show the global shared-memory limits and utilization summary
\end{lstlisting}

\begin{notebox}
ToSh wraps `lsipc -J` so IPC resources and global IPC limits flow through the pipeline as structured rows instead of terminal-shaped text. Text-format-only modes like `--raw`, `--export`, `--newline`, and shell-variable output currently fall back to the external `lsipc` utility unchanged.
\end{notebox}

\end{cmdbox}

\subsection{\texttt{networkctl}}
\label{ref:networkctl}
\icmd{networkctl}

\begin{cmdbox}{networkctl}
Wraps networkctl, returning typed network-link rows for supported list queries.

\begin{signature}
networkctl [list [pattern ...] | <other-command ...>]
\end{signature}

\begin{tabular}{@{}>{\ttfamily\small\raggedright\arraybackslash}p{5cm} L{9cm}@{}}
\toprule
\textbf{Argument} & \textbf{Description} \\
\midrule
{}[list [pattern ...]] & With no subcommand, ToSh treats `networkctl` as a structured `list` query. Explicit `list` behaves the same way. (optional) \\
<other-command ...> & Unsupported, detail, and mutating commands currently fall back to the native `networkctl` utility unchanged. (optional) \\
\bottomrule
\end{tabular}

\begin{tabular}{@{}>{\ttfamily\small\raggedright\arraybackslash}p{5cm} L{9cm}@{}}
\toprule
\textbf{Flag / Option} & \textbf{Description} \\
\midrule
-a|\allowbreak --all & Pass through to the structured `networkctl list` query to include all visible links. \\
-l|\allowbreak --full & Pass through to the structured `networkctl list` query. \\
--show <columns> & Use ToSh display-only column selection on the structured network-link rows. \\
--hide <columns> & Hide display columns while preserving the underlying typed link objects. \\
--show-all & Expose every selectable structured network-link display column. \\
\bottomrule
\end{tabular}

\textbf{Pipeline input:} none
 --- The structured `networkctl` builtin is explicit-arg-first and does not currently consume pipeline input.

\textbf{Output:} Returns typed network-link rows for supported `networkctl list` queries. Other commands currently fall back to the native `networkctl` output.

\begin{lstlisting}
networkctl  # List network links as typed rows
networkctl | where _.Setup == unmanaged  # Filter links by setup state in the pipeline
networkctl --show Link,Operational,Setup,Managed  # Render a focused network-link summary
\end{lstlisting}

\end{cmdbox}

\subsection{\texttt{seq}}
\label{ref:seq}
\icmd{seq}

\begin{cmdbox}{seq}
Generates a numeric sequence.

\begin{signature}
seq <stop> | seq <start> <stop> | seq <start> <step> <stop>
\end{signature}

\end{cmdbox}

\subsection{\texttt{sleep}}
\label{ref:sleep}
\icmd{sleep}

\begin{cmdbox}{sleep}
Pauses execution for a duration.

\begin{signature}
sleep <duration> [duration...]
\end{signature}

\end{cmdbox}

\subsection{\texttt{systemctl}}
\label{ref:systemctl}
\icmd{systemctl}

\begin{cmdbox}{systemctl}
Wraps systemctl, returning typed unit rows for supported list queries and structured unit property sets for supported show and status queries.

\begin{signature}
systemctl [list-units [pattern ...] | list-unit-files [pattern ...] | \\
\quad show <unit ...> [options] | status <unit ...> | <other-subcommand \\
\quad ...>]
\end{signature}

\begin{tabular}{@{}>{\ttfamily\small\raggedright\arraybackslash}p{5cm} L{9cm}@{}}
\toprule
\textbf{Argument} & \textbf{Description} \\
\midrule
{}[list-units [pattern ...]] & With no subcommand, ToSh treats `systemctl` as a structured `list-units` query. Explicit `list-units` behaves the same way. (optional) \\
list-unit-files [pattern ...] & Returns typed unit-file state rows for installed unit files. (optional) \\
show <unit ...> & Returns structured systemd unit property sets for one or more units. (optional) \\
status <unit ...> & Returns structured unit status objects for one or more units, including recent logs when available. (optional) \\
<other-subcommand ...> & Unsupported subcommands fall back to the native `systemctl` utility unchanged. (optional) \\
\bottomrule
\end{tabular}

\begin{tabular}{@{}>{\ttfamily\small\raggedright\arraybackslash}p{5cm} L{9cm}@{}}
\toprule
\textbf{Flag / Option} & \textbf{Description} \\
\midrule
--type <type[,type...]>|\allowbreak -t <type[,type...]> & Restrict structured `list-units` or `list-unit-files` output to specific unit types such as `service` or `socket`. \\
--state <state[,state...]> & Restrict structured `list-units` output to specific load/active states. \\
--all & Include inactive and unloaded units in the structured listing. \\
--failed & Restrict the structured listing to failed units. \\
-p <property[,property...]>|\allowbreak --property <property[,property...]> & Restrict `show` to specific fetched properties. ToSh still injects `Id` internally so multiple-unit output stays structured. \\
--show <columns> & Use ToSh display-only column selection on structured unit rows or property sets. \\
--hide <columns> & Hide display columns while preserving the underlying typed objects. \\
--show-all & Expose every selectable structured display column for the current output shape. \\
\bottomrule
\end{tabular}

\textbf{Pipeline input:} none
 --- The structured `systemctl` builtin is explicit-arg-first and does not currently consume pipeline input.

\textbf{Output:} Returns typed systemd unit rows for `list-units`, typed unit-file rows for `list-unit-files`, and structured unit property-set objects for supported `show` and `status` queries. Other subcommands currently fall back to the native `systemctl` output.

\begin{lstlisting}
systemctl  # List units as typed rows
systemctl --type service | where _.Active == active  # Filter active services in the pipeline
systemctl list-unit-files --type service | where _.Enabled  # Inspect installed enabled service unit files
systemctl status sshd.service | get { Id, ActiveState, MainPID, RecentLogCount }  # Inspect structured unit status details
\end{lstlisting}

\end{cmdbox}

\subsection{\texttt{timespan}}
\label{ref:timespan}
\icmd{timespan}

\begin{cmdbox}{timespan}
Parses a duration into a CLR duration value.

\begin{signature}
timespan <duration>
\end{signature}

\end{cmdbox}

\subsection{\texttt{uname}}
\label{ref:uname}
\icmd{uname}

\begin{cmdbox}{uname}
Returns kernel and operating system information as a Tosh object.

\begin{signature}
uname [-a|-s|-n|-r|-v|-m|-o]
\end{signature}

\begin{tabular}{@{}>{\ttfamily\small\raggedright\arraybackslash}p{5cm} L{9cm}@{}}
\toprule
\textbf{Flag / Option} & \textbf{Description} \\
\midrule
-s & System Name \\
-n & Node Name \\
-r & Kernel Release \\
-v & Kernel Version \\
-m & Machine \\
-o & Operating System \\
\bottomrule
\end{tabular}

\end{cmdbox}

\subsection{\texttt{unset}}
\label{ref:unset}
\icmd{unset}

\begin{cmdbox}{unset}
Removes Tosh variables, functions, and exported environment names.

\textit{Aliases:} \code{forget}

\begin{signature}
unset <name> [name...]
\end{signature}

\end{cmdbox}

\subsection{\texttt{uptime}}
\label{ref:uptime}
\icmd{uptime}

\begin{cmdbox}{uptime}
Returns system uptime and load averages.

\begin{signature}
uptime
\end{signature}

\end{cmdbox}

\subsection{\texttt{vars}}
\label{ref:vars}
\icmd{vars}

\begin{cmdbox}{vars}
Lists all visible variables in the current scope.

\begin{signature}
vars [filter]
\end{signature}

\end{cmdbox}

\subsection{\texttt{whoami}}
\label{ref:whoami}
\icmd{whoami}

\begin{cmdbox}{whoami}
Returns the current user principal.

\begin{signature}
whoami
\end{signature}

\end{cmdbox}


\section{Text}

\subsection{\texttt{cut}}
\label{ref:cut}
\icmd{cut}

\begin{cmdbox}{cut}
Extracts character or delimited fields from text.

\begin{signature}
cut (-f fields [-d delimiter] | -c chars) [path ...]
\end{signature}

\end{cmdbox}

\subsection{\texttt{echo}}
\label{ref:echo}
\icmd{echo}

\begin{cmdbox}{echo}
Emits its arguments as pipeline objects.

\begin{signature}
echo <value> [value...]
\end{signature}

\begin{tabular}{@{}>{\ttfamily\small\raggedright\arraybackslash}p{5cm} L{9cm}@{}}
\toprule
\textbf{Argument} & \textbf{Description} \\
\midrule
value & One or more values to emit as pipeline objects. \\
\bottomrule
\end{tabular}

\textbf{Output:} Each argument as a separate pipeline object.

\begin{lstlisting}
echo hello world
echo 42 | where { $_ > 10 }
\end{lstlisting}

\end{cmdbox}

\subsection{\texttt{grep}}
\label{ref:grep}
\icmd{grep}

\begin{cmdbox}{grep}
Searches text input with a regular expression or literal pattern.

\begin{signature}
grep [-i] [-v] [-F] [-n] [-r] [-w] [-o] [-c] [-l] [-L] [-A num] [-B num] \\
\quad [-C num] <pattern|regex> [path ...]
\end{signature}

\begin{tabular}{@{}>{\ttfamily\small\raggedright\arraybackslash}p{5cm} L{9cm}@{}}
\toprule
\textbf{Argument} & \textbf{Description} \\
\midrule
pattern|\allowbreak regex & The text pattern or .NET regular expression to search for. \textit{(string|regex)} \\
path & Optional file paths to search instead of consuming piped text. (optional) \textit{(path-like)} \\
\bottomrule
\end{tabular}

\begin{tabular}{@{}>{\ttfamily\small\raggedright\arraybackslash}p{5cm} L{9cm}@{}}
\toprule
\textbf{Flag / Option} & \textbf{Description} \\
\midrule
-i & Ignore case. \\
-v & Invert the match. \\
-F & Treat the pattern as a literal string instead of a regex. \\
-n & Include source line numbers in text-file results. \\
-m & Multiline mode. \\
-s & Singleline mode so '.' matches newlines. \\
-x & Require a full-line match. \\
--explicit-capture & Return structured capture results instead of plain matching lines. \\
\bottomrule
\end{tabular}

\textbf{Pipeline input:} scalar
 --- Consumes scalar text from the pipeline. When paths are supplied explicitly, grep reads file contents instead.

\textbf{Output:} Matching text lines, or structured regex capture objects with --explicit-capture.

\begin{lstlisting}
echo one two three | grep tw  # Pipe text into grep
echo "Alpha" | grep -i "^alpha$"  # Use regex flags
grep -F literal README.md  # Search a file literally
\end{lstlisting}

\end{cmdbox}

\subsection{\texttt{head}}
\label{ref:head}
\icmd{head}

\begin{cmdbox}{head}
Returns the first objects or first text lines.

\begin{signature}
head [-n count] [-c bytes] [path ...]
\end{signature}

\end{cmdbox}

\subsection{\texttt{join-lines}}
\label{ref:join-lines}
\icmd{join-lines}

\begin{cmdbox}{join-lines}
Joins input values into a single text value.

\begin{signature}
join-lines [separator]
\end{signature}

\end{cmdbox}

\subsection{\texttt{lines}}
\label{ref:lines}
\icmd{lines}

\begin{cmdbox}{lines}
Splits text input into individual lines.

\begin{signature}
lines [text ...]
\end{signature}

\end{cmdbox}

\subsection{\texttt{match}}
\label{ref:match}
\icmd{match}

\begin{cmdbox}{match}
Matches text with a regular expression and returns shell record objects.

\begin{signature}
match [-i] [-m] [-s] [-x] [--explicit-capture] <pattern|regex> [text \\
\quad ...]
\end{signature}

\begin{lstlisting}
echo "PID=42" | match "PID=(?<Pid>[0-9]+)" | get Pid
echo "Alpha" | match -i "^alpha$"
\end{lstlisting}

\end{cmdbox}

\subsection{\texttt{raw}}
\label{ref:raw}
\icmd{raw}

\begin{cmdbox}{raw}
Emits plain text without rich table or record rendering.

\begin{signature}
raw [value...]
\end{signature}

\begin{tabular}{@{}>{\ttfamily\small\raggedright\arraybackslash}p{5cm} L{9cm}@{}}
\toprule
\textbf{Argument} & \textbf{Description} \\
\midrule
value ... & Optional explicit values to emit as plain text instead of rich object display. (optional) \\
\bottomrule
\end{tabular}

\textbf{Pipeline input:} scalar, record, list, table
 --- Consumes pipeline values and emits one plain-text line per input item.

\textbf{Output:} Returns ShellTextLine values so the final display stays plain text instead of tables or record views.

\begin{lstlisting}
echo 1317 | raw  # Show a scalar without rich boxing
echo System.DayOfWeek.Friday | raw  # Show an enum's raw CLR string form
\end{lstlisting}

\end{cmdbox}

\subsection{\texttt{replace}}
\label{ref:replace}
\icmd{replace}

\begin{cmdbox}{replace}
Replaces text in each input value.

\begin{signature}
replace [-i] [-m] [-s] [-x] [--explicit-capture] [-r] <pattern|regex> \\
\quad <replacement> [text ...]
\end{signature}

\begin{lstlisting}
echo alpha-beta | replace beta BETA
echo "A1 B2" | replace -r "[0-9]" "#"
\end{lstlisting}

\end{cmdbox}

\subsection{\texttt{split}}
\label{ref:split}
\icmd{split}

\begin{cmdbox}{split}
Splits text values into smaller text values.

\begin{signature}
split [-r] [-i] [-m] [-s] [-x] [--explicit-capture] [delimiter|regex] \\
\quad [text ...]
\end{signature}

\begin{lstlisting}
echo "alpha,beta,gamma" | split ","
echo "alpha,beta;gamma" | split -r "[,;]"
\end{lstlisting}

\end{cmdbox}

\subsection{\texttt{tail}}
\label{ref:tail}
\icmd{tail}

\begin{cmdbox}{tail}
Returns the last objects or last text lines.

\begin{signature}
tail [-n count] [-c bytes] [-f] [path ...]
\end{signature}

\end{cmdbox}

\subsection{\texttt{template}}
\label{ref:template}
\icmd{template}

\begin{cmdbox}{template}
Renders pipeline objects into text using \{\{ member.path \}\} placeholders.

\begin{signature}
template <text>
\end{signature}

\end{cmdbox}

\subsection{\texttt{tr}}
\label{ref:tr}
\icmd{tr}

\begin{cmdbox}{tr}
Translates or deletes characters in text.

\begin{signature}
tr [-d] <set1> [set2] [path ...]
\end{signature}

\end{cmdbox}

\subsection{\texttt{uniq}}
\label{ref:uniq}
\icmd{uniq}

\begin{cmdbox}{uniq}
Collapses adjacent duplicate input values.

\begin{signature}
uniq [-c] [-i] [path ...]
\end{signature}

\end{cmdbox}

\subsection{\texttt{wc}}
\label{ref:wc}
\icmd{wc}

\begin{cmdbox}{wc}
Counts lines, words, bytes, characters, and longest-line length.

\begin{signature}
wc [-lwmcL] [path ...]
\end{signature}

\begin{tabular}{@{}>{\ttfamily\small\raggedright\arraybackslash}p{5cm} L{9cm}@{}}
\toprule
\textbf{Argument} & \textbf{Description} \\
\midrule
path ... & Optional file paths to measure. When omitted, `wc` measures the current text pipeline. (optional) \textit{(path-like)} \\
\bottomrule
\end{tabular}

\begin{tabular}{@{}>{\ttfamily\small\raggedright\arraybackslash}p{5cm} L{9cm}@{}}
\toprule
\textbf{Flag / Option} & \textbf{Description} \\
\midrule
-l & Show line counts. \\
-w & Show word counts. \\
-c & Show byte counts. \\
-m & Show character counts. \\
-L & Show the longest-line length. \\
\bottomrule
\end{tabular}

\textbf{Pipeline input:} scalar, record
 --- Consumes the current text pipeline when explicit file paths are omitted.

\textbf{Output:} Returns typed text-statistics objects, and appends a `total` row when multiple files are counted.

\begin{lstlisting}
wc README.md
wc -lwm README.md
echo one two three | wc
\end{lstlisting}

\begin{notebox}
Wc returns typed statistics objects instead of formatted text, so you can still `get`, `where`, or `summarize` them later. Selector flags like `-l` and `-w` only change the visible columns, not the underlying objects.
\end{notebox}

\end{cmdbox}

\subsection{\texttt{write}}
\label{ref:write}
\icmd{write}

\begin{cmdbox}{write}
Writes rendered values without a trailing newline.

\begin{signature}
write [value...]
\end{signature}

\end{cmdbox}

\subsection{\texttt{writeline}}
\label{ref:writeline}
\icmd{writeline}

\begin{cmdbox}{writeline}
Writes rendered values with a trailing newline.

\begin{signature}
writeline [value...]
\end{signature}

\end{cmdbox}

